\section{The sources block}\label{sec:sources}

This section gives one paragraph for each project file \file{SOURCES_*.md} that has a verifier file. Each paragraph states the result, describes the method in two sentences, gives the verifier's grade in the tags of Section~\ref{sec:status}, and names the inputs that stay open. Nothing here goes beyond the project ledgers: the verifier files, \file{TRUST_LEDGER.md}, the dependency table (Section~\ref{sec:deptable}), and \file{PLAN.md}. Where a verifier corrected a file, the corrected statement is given. The grouping follows \S4 of \file{TRUST_LEDGER.md}: first the files that Theorem 2w (Target 1) uses, then the files of the finite-torus theorem (Target 2) in the order of its dependency chain (steps 2--9 of \S4.2 there). Ledger labels A$n$, B$n$, C$n$, X$n$ are those of \file{TRUST_LEDGER.md}. \file{SOURCES_A1_KERNEL.md} and \file{SOURCES_A4_D189.md} share the verifier \file{SOURCES_A1_A4_VERIFY.md}, and \file{SOURCES_HHIST.md} is checked in \file{SOURCES_HHIST_N1_VERIFY.md}. Three files have no verifier and are not listed: \file{SOURCES_BOOKKEEPING.md}, \file{SOURCES_G4.md}, \file{SOURCES_W4.md}. Files verified after this section was assembled have no paragraph of their own: \file{SOURCES_WC_EXPGAUGE.md} and \file{SOURCES_WC_RESIDUE.md} (cited under A1w and A2), and \file{SOURCES_T2_V1.md}, \file{SOURCES_T2_PI.md}, \file{SOURCES_T2_A1WO3.md}, \file{SOURCES_T2_A1WO4.md}, \file{SOURCES_T2_D189.md}, \file{SOURCES_T2_R2REM.md}; their state is in \file{TRUST_LEDGER.md} and \file{PLAN.md}.

\subsection{Target 1: inputs of Theorem 2w}\label{sec:src-T1}

\paragraph{\file{SOURCES_GGAIN_NONMONO.md}: item 7$'$, the running-maximum lift.}
For every step $k$, every past $h\in\ioc{\gamma}^k$, monotone or not, and every $g\in[0,\gamma]$, $|\beta_{k+1}(h,g)-\beta_{k+1}(0)|\le C_\beta g$ with $C_\beta=C_\Pi S_0$, the constant of (G-gain), uniform in $k$, $K$, and the past. This removes the monotone-history restriction of (G-gain), which blocked item 7$'$ on the clamped, non-monotone pasts that the shooting evaluates.
The only use of monotonicity is Step 2 of Lemma F2$'$, where $\max_{j\le k}4g_j^{\kappa_0}=4g_k^{\kappa_0}$ compares bound sizes; with the running maximum $G_k=\max_{i\le k}g_i$ the induction $a_j\le2G_j^{\kappa_0}$ closes for every past under the side condition $4\gamma^{\kappa_0}\le E_0$, which is of Ba{\l}aban's existing type. The value at zero own coupling, (K1), is history-free because the exponent vanishes at $g_k=0$ (CMP 109 p.\ 268) for any input, and the endpoint $g\to0$ is a dominated-convergence step that is not written line by line.
Verifier: \file{SOURCES_GGAIN_NONMONO_VERIFY.md}, the lift \status{proved+verified}, item 7$'$ \status{proved+verified, conditional on (H0) (ledger A5) or on the (G-gain) inputs (ledger B13)}; all four items confirmed with small corrections: $C_\beta'=C_\beta$ holds because $E_0+4\gamma^{\kappa_0}\le2E_0$, the input radius that (G-gain) already covers, and the possible dependence of the newest R-input on $g$ (caveat c1) is open on the literal reading of $\beta$ and void on the pure-part reading.
Open inputs: (H0) (ledger A5; its use here is a verified reduction with residue N1--N9, Section~\ref{sec:src-r5}), or the (G-gain) inputs of ledger B13 ((Q)/$c_-$, ($M_r$), CMP 102 (189), the (M-B) bookkeeping); L5, the evaluation of $\beta$ off Ba{\l}aban's trajectory (item 8, ledger B7); caveat c4, CMP 119 (2.31) and the off-trajectory small-field induction of \file{SOURCES_SM.md} \S5.1 not swept for coupling-order uses.

\paragraph{\file{SOURCES_T2_KUNIF.md}: item 9, $K$-independence of $\beta_j(g)$.}
Item 9 is a volume statement, not a history identity, and it closes as a reduction: CMP 119 (2.24), p.\ 253, and p.\ 279 define $\beta_j$ as the whole-lattice coefficient of CMP 109 (whose $\Pi$ in (4.37) is a sum over $\mathbb Z^4$), so the $\beta_j$ of the recursion is $K$-free by definition. The shooting needs only (K1) $\beta_j(h,0)=\beta_j(0)$ for every admissible past, (K2) $|\partial\beta_j/\partial g_{j-1}|\le C_\beta$, and (K3) the window bound $\sum_{i=k+1}^{n}\beta_i(0)\in[b(n-k)-2R,\,b(n-k)+2R]$ of Theorem K; with $\gamma\le b/(4C_\beta)$ these give $T(m,n)\in[(3b/4)(n-m)-2R,\,(5b/4)(n-m)+2R]$ on every admissible history.
Lemmas D$'$ and D$''$ run at one fixed $K$ and never compare two trajectories or two values of $K$, so only their constants must be free of $K$. The bound $|\beta_{k+1}(h_k)-\beta_{k+1}(0)|\le C_\beta g_k$ follows from (K1) and the mean-value theorem in the own slot, and it gives the window bracket pointwise on each clamped history of the intermediate-value argument.
Verifier: \file{SOURCES_T2_KUNIF_VERIFY.md}, \status{verified reduction}; items 1, 2, 4 confirmed, item 3 corrected: (G-gain) gave $C_\beta$ only on monotone slot histories, so (K2) via (G-gain) needs item 7$'$ (now lifted by \file{SOURCES_GGAIN_NONMONO.md}); the CMP 119 page for the definition of $\beta_j$ is p.\ 253, not p.\ 249.
Open inputs: item 7$'$ for (K2) (from (H0), ledger A5, or from (G-gain), ledger B13); L5, the domain and continuity of $\beta$ off Ba{\l}aban's trajectory (item 8, ledger B7); the inputs of Theorem K for (K3); the finite-volume identification of the step (ledger B8, with A15, narrowed by Proposition~\ref{prop:A15N}), an input of Theorem 1, not of the shooting.

\paragraph{\file{SOURCES_T2_V1.md}: item 5, V1, ledger B5.}
Inside Ba{\l}aban's construction with the cut CMP 109 (2.9), the term $E^{(k+1)}$ of (2.13) vanishes identically at $g_k=0$. So $\Pi^{(k+1)}(0)$ and $\beta_{k+1}(0)$ come from $\log Z^{(k)}$ alone: $\beta_{k+1}(0)$ is the Maxwell coefficient of the one-step Gaussian determinant, it is free of the history, and it equals the determinant of the numerics up to a constant.
The file reads the definitions (1.3)--(1.22) and (2.10)--(2.15) on page images and passes from Ba{\l}aban's exponential gauge fixing to the bordered Hessian of the numerics (gauge-fixing independence, constraint Jacobian, expansion point, normalization).
Verifier: \file{SOURCES_T2_V1_VERIFY.md}, \status{proved+verified, conditional} at fixed level, with 14 claims confirmed and 3 corrected; the conditions are (G2), Ba{\l}aban's own inductive bounds, and a declaration that the averaging is (0.4) (the numerics use (0.4); Ba{\l}aban's step uses (0.12)).
Open inputs: (G2) (item 6); the averaging declaration; Ba{\l}aban's inductive bounds.

\paragraph{\file{SOURCES_T2_G2.md}: item 6, (G2), ledger B6.}
At fixed level $j$, the constrained one-loop log-determinant on a torus of side $N$ has no zero modes, and its density per coarse cell converges to the infinite-lattice density with error $\le C_je^{-\delta_jN}$, also for every derivative in the plane-wave amplitude. The $p^2$-coefficient commutes with $N\to\infty$; the $f^2$-coefficient at the constant field commutes with $N\to\infty$, given a gap input (Gap$_f$) and $C^4$-regularity in $f$, both [R]. The massless pieces of Prop.~2.1 are handled fiberwise, with an $N$-independent zero-mode constant.
The proof uses a Combes--Thomas decay bound and the method of images; the constants are fixed-$j$ constants ($\delta_j\gtrsim L^{-2j}$), and the $j$-uniform version is not proved.
Verifier: \file{SOURCES_T2_G2_VERIFY.md}, 21 claims confirmed and 6 corrected, none affecting a conclusion; item 6 moves from \status{formal} to [R] (\status{verified reduction}) at fixed $j$.
Open inputs: the surjectivity of $\tilde Q$, (Gap$_f$), and the reading that the averaging is analytic.

\paragraph{\file{SOURCES_T2_S123.md}: ($M_j$), ledger B4.}
At fixed $j\ge1$, odd $L\ge3$, and small flux $f$, the lattice spectral inputs hold: (S1) $\mathcal K_t(f)\ge-C_1|f|$ for $t\in[0,1]$, (S3) $\pm V\le C_3[(\mathcal K_t+\mu)^2+f^2]$ in the subdivided form of \file{KAPPA_BOUNDED.md} \S1.3, (S2) at $t=0$, and the coarse comparison (S2c) $N_0(\lambda-\eta_1)\le N_t(\lambda)\le N_0(\lambda+\eta_1)$ with $\eta_1=O(f^2)$. Corollary 7.1 then reduces ($M_j$), the limit of $[\Gamma_c(\Lambda_j;f)-\Gamma_c(1;f)-(f=0\text{ value})]/(n_2f^2)$ to $d_j$, to the near-zero input $(Z_j)$, itself reduced to the 2D gap (G$_A'$), and to region 1 of Prop.\ 4.4 of \file{BETA_FP_M.md}.
The proofs work in the twisted algebra of operators that commute with the magnetic translations, with Peierls products, Moyal expansions with kernel-norm remainders, a Nenciu-type inverse, and operator bounds from the vanishing order of symbols. At $t=0$ a fibre decomposition over $k_\perp$ separates the operator, and a gauge-copy lemma (Lemma 5.8) shows that the near-zero density of states needs only (G$_A'$): the transverse charged 2D SU(2) Wilson Hessian at constant flux has no spectrum in $(-c|f|,c|f|)$, and numerically its gap is about $|ef|$.
Verifier: \file{SOURCES_T2_S123_VERIFY.md}, ($M_j$) \status{verified reduction}, with (S1), (S3), (S2)$|_{t=0}$, (S2c) and region 2 of Prop.\ 4.4 \status{proved+verified at fixed $j$, conditional on (I3)}; 25 claims confirmed, 2 corrected (in Lemma 1.4 $\nabla^\beta$ is one twisted convolution; the explicit $L$-power of $C_j$ in Prop.\ 3.2 is convention-dependent and is not quoted), and the input list of ($M_j$) gains the neutral-sector differentiation under the integral, an inference.
Open inputs: (I3), the entrywise analyticity in $f$ of Ba{\l}aban's linearized averaging at each of the $j$ steps, is now proved (next paragraph); (G$_A'$) (open, independent of $j$ and $L$, numerically confirmed at six rational fluxes with margin about $|ef|$); region 1 of Prop.\ 4.4 of \file{BETA_FP_M.md} (argued); all under ledger B4.

\paragraph{\file{SOURCES_T2_I3.md}: the Peierls input (I3), ledger B4.}
Let $L$ be odd $\ge3$, and let the averaging be Ba{\l}aban's (0.4), or (0.12) with Federbush's average. Then the linearized averaging $\tilde Q_f$ at the constant-flux background is entrywise analytic in complex $f$ on $|f|<f_0$, $f_0=\pi/8$ for (0.4) and $\min(\pi/8,1/8)$ for (0.12), and (I3) holds at every step $i\le j$ with radius and constants independent of $i$ and $j$ (Theorem 7.1 there). It holds in the stated two-block range and does not use (I1) for the fixed vertices.
The proof linearizes the averaging explicitly at the flux background (Prop.\ 5.3), bounds each charged entry of $\tilde Q_f$ and of $D_f$ by its reduced kernel (Props.\ 5.4, 6.1), and assembles the bordered matrix of one step.
Verifier: \file{SOURCES_T2_I3_VERIFY.md}, 15 claims confirmed, 2 corrected; (I3) moves from [I] to [P], given the identification of the one-step objects of the project with the linearization of (0.4) or (0.12) at the constant-flux background.
Open inputs: that identification; for ($M_j$) the 2D gap (G$_A'$) and region 1 of Prop.\ 4.4 of \file{BETA_FP_M.md} stay.

\subsection{Target 2, step 2: the sourced Theorem 1 (ledger B11)}\label{sec:src-B11}

\paragraph{\file{SOURCES.md}: the E14 reduction of (B-S).}
The plaquette-density source is a local bare coupling, $\rho_0^f=\exp(-\sum_p(g_0^{-2}-f_p)P_p(U))$, so the Schwinger functions on test functions flat on the diagonals reduce to a sourced version of Ba{\l}aban's Theorem 1 with a pointwise coupling flow (S-A), sourced regular terms (S-B), a local contact functional $W_k[f]$ (S-C), a local $\beta$-function (S-D), and the normalization (S-E). The new estimates are (E14-a) the sourced small-field step, (E14-b) the sourced large-field step and R operation, and (E14-c) the derivative bounds in $f$; (E14-d) = E13 is needed only for full-sequence convergence, and (E14-f) only for nontriviality.
A dimension count shows that the source of the dimension-4 density is marginal at every degree, that its field-independent part is relevant and is removed as local contact terms, and that the first field-dependent gradient term has dimension 6. One step is written out at the Gaussian level: with $K_f=(\Delta^{(1)})^{-1/2}\Delta^{(f)}(\Delta^{(1)})^{-1/2}$ and $\|K_f\|\le r$ the source enters through $-\tfrac12\log\det(1-g_0^2K_f)$, and the linear response per step is $c_j=1+\tfrac12g_{j-1}^3\beta_j'(g_{j-1})$.
Verifier: \file{SOURCES_VERIFY.md}, \status{verified reduction} (the structure confirmed; items 5, 6, 8 corrected, item 10 confirmed with corrections); the printed CMP 119 (2.24) is $1/g_{j-1}^2(x)=1/g_j^2(x)+\beta_j(g_{j-1})\varphi_j(x)$, and by (2.25) Ba{\l}aban evaluates $\beta_j$ at the constant coupling, so (S-A) and (S-D) are new; derivatives of $\log Z$ are cumulants; $\prod_jc_j$ is uniformly bounded, not shown convergent, and assumes a one-step recursion; the uniform bound for general $F$ flat on the diagonals is not derived from the source form.
Open inputs: (B-S) as a whole, (E14-a, b, c) (ledger B11, \status{open}), with its parts B12--B20 and B22; the (E14-c) clause of the bridge (B-N1) (ledger B21); (NT) = (E14-f) for Theorem~\ref{thm:nontriv}; E13 = (E14-d), not needed for Theorem~\ref{thm:B}.

\paragraph{\file{SOURCES_D1.md}: one sourced step at one loop (D1).}
At one loop the plaquette-density source, a local bare coupling $w_p=1/g_0^2-f_p$, gives the multiplicative source renormalization $c_1=1$ exactly, for every background; the first correction is $c_1-1=\tfrac12g^3\beta_1'(g)$, which is $\beta_1^{(2)}g_0^4+O(g_0^6)$ if $\beta_1$ is even and $O(g_0^3)$ with a $C^1$ bound only. The field-independent terms are the tadpole $\tfrac12(9L^4-9)g_0^2$ per coarse cell ($\tfrac34(1-L^{-4})g_0^2$ per plaquette) and a contact term $C(q)$ that is analytic at $q=0^+$, and the first field-dependent gradient terms have dimension 6.
The source enters the step only through the Hessian $\Delta_f=\sum_pw_pH_p$ and a current term $\langle J_f,B\rangle$ at Ba{\l}aban's unweighted minimizer, while gauge fixing, ghosts, and the constraint carry no $f$; for constant $f$ the quadratic form scales by $1-g_0^2f_0$, so the determinant changes by a field-independent constant over $9L^4-9$ modes per cell. Per-momentum formulas for the tadpole, $C(q)$, and the mixed derivative are evaluated by \file{numerics/d1_check.py} and checked against real-space finite differences on the periodic $6^4$ lattice.
Verifier: \file{SOURCES_D1_VERIFY.md}, \status{proved+verified} for the algebra (items 1--4, 6, 8--10 confirmed) and \status{numerical, reproduced} for the rerun, which matched every digit; corrected: the dimension-6 count holds by index parity, the $\gamma_2$ values are unreliable, the data support $|c_1-1|\le5\times10^{-4}$ but not geometric convergence in the grid, and the expansion-point error has the small parameter $Lp$.
Open inputs: none at the Gaussian level; the non-Gaussian sourced step and the pointwise flow belong to D2--D3 and to (B-S), ledger B11.

\paragraph{\file{SOURCES_D2.md}: sourced small-field step (E14-a).}
The sourced small-field step is Ba{\l}aban's step with extra source vertices of order $g_k$, $g_k^2$, $g_k^3$, all in his $P^{(k)}$ class (each $\le O(g_k)(1+|B|^2)$ on $|B|<\varepsilon_1/g_k$), a propagator for the sourced Hessian with Ba{\l}aban's decay and regularity bounds (constant $3c_1$ with $c_1=O(\delta_0^{-5})c_0$, rate $\delta_0/2$), and a field-independent vacuum term $W$ that factors out. Its step (c), the pointwise flow $1/g_{j+1}^2(x)=1/g_j^2(x)-B_{j+1}(f(x))$ with $h\equiv0$, is not established here and is redone in \file{SOURCES_D2C.md}.
The source term is evaluated at Ba{\l}aban's unweighted minimizer, so his measure and propagators are used without change, and the sourced propagator follows from a weighted Neumann series in his cube-to-cube operator norms (CMP 116 (1.7)). Real-source derivative bounds are tied to polystrip holomorphy by the interpolation Lemma 7.1, and the remainder bounds are reduced to Ba{\l}aban's estimates with a freezing property and a constant-source family.
Verifier: \file{SOURCES_D2_VERIFY.md}, steps (a), (b), (d), (f) \status{proved+verified}, step (e) a \status{verified reduction}, and step (c) rejected, since its claim that CMP 109 Sects.~3--4 run verbatim is \status{refuted} ((4.34)--(4.37) use Euclidean symmetry, and the cubes are at scale $k$, not $j$); corrected: (a) needs a block-constant $\bar f$, the constant is $3c_1$ and not $\times3$, $W$ factors only in the small-field region, and the polydisc radius is $r/(2n)$.
Open inputs: analyticity of the step in its input activities, (W4), and the field-dependent domains $\Lambda_j$ in $W$ (ledger B20); the recovery $\kappa\to\kappa/2$ (ledger B22, a verified reduction that uses A23); (H$\beta$), reduced in \file{SOURCES_HBETA.md}; the scale condition (U) is the constant choice C23, not an open input.

\paragraph{\file{SOURCES_D2C.md}: step (c) redone, frozen source per domain (Lemma A).}
Each step term $E^{(j)}(X,\cdot\,;f)$, $X\in D_j$, is frozen at $c_X=f(b_X)$, and the freezing error obeys $|D_X|\le E_0CM^3(\rho_1/r)(L^j\eta)^5\lambda_ke^{-\kappa d/2}$ (Lemma A), irrelevant with exponent 1 and with no growth in $k$. With the identification (I) of the constant-source increment $B_j(c)$ with Ba{\l}aban's $\beta_j$ on the trajectory from $1/g_0^2-c$, this gives the pointwise recursion $1/g_{j+1}^2(x)=1/g_j^2(x)-B_{j+1}(f(x))$ with $h\equiv0$ and gradient corrections irrelevant with exponent $\min(\alpha,\beta,1)$.
The non-constant source is handled at scale $j$, per localization domain, as a gauge-invariant difference bounded in an axial gauge on $X$, where each field costs two powers of $L^j\eta$. Ba{\l}aban's CMP 109 Sects.~4--5 are then used only for constant sources, where they apply at bare coupling $1/g_0^2-c$.
Verifier: \file{SOURCES_D2C_VERIFY.md}, Lemma A and the recursion \status{proved+verified, conditional on (D) and (H$\beta$)}; corrected: the freezing points are the points of $X$ nearest to $b_X$, (I) is exact only with the adapted gauge-fixing weight $1/u_k(c)$ (otherwise up to $O(e^{-c\varepsilon_1^2/g_k^2})$), $\beta>0$ in (4.18) must be fixed, and the exact cancellation of the cubic first-moment term of (4.34) is unsupported and not needed.
Open inputs: the domain hypothesis (D) (ledger B20; argued in \file{SOURCES_D2C_VERIFY.md} \S2 from the CMP 109 p.~263 remark, which cites CMP 102 Prop.~9, ledger A4); (H$\beta$) (Lemma H of \file{SOURCES_HBETA.md}); (I) with large-field thresholds (ledger B12); the adapted gauge-fixing weight changes the construction (ledger C10).

\paragraph{\file{SOURCES_D3.md}: sourced large-field step (E14-b).}
The sourced large-field step is Ba{\l}aban's step with all thresholds frozen at the reference couplings, the local coupling $w_j=w_{j-1}-B_j(f)\varphi_j$, an R operation that resets it on $Z$ to $u_k(f(x))$ (CMP 122 II p.~365), and an adapted local gauge-fixing weight. Real sources keep the positivity of the large-field weights under $3\gamma^2r\le\tfrac12$ (revised to $5\gamma^2r\le\tfrac12$ in \file{SOURCES_D3R.md}), Lemma D3.1 gives the local holomorphy radius $c_\tau r/(E_0|Y|)$ of one positive denominator, and the identification (I) is off by a threshold term, the new input (H-thr).
The file lists the places (P1)--(P7) where Ba{\l}aban uses a constant coupling in the large-field step and states the change at each. Each change is reduced to a quoted Ba{\l}aban estimate (CMP 119 (2.23)--(2.31), CMP 122 II (1.73)--(1.75)); a line-by-line rerun of CMP 122 I--II is not done.
Verifier: \file{SOURCES_D3_VERIFY.md}, \status{verified reduction} with corrections: the reset term is bounded by Ba{\l}aban's sup bound (CMP 122 II p.~366), not cancelled; the t-parameter term (P8) of CMP 119 p.~271 is missing; old R-terms are merged into $E_k$, so freezing through R is open again; the Gevrey order is $(n!)^{2+\epsilon}$, not $(n!)^2$; and the lower bound on $w_k$ holds relative to Ba{\l}aban's local coupling.
Open inputs: (H-thr) (ledger B15); the real uniform (H-hist) (ledger B13); strip holomorphy of $u_k$, now N3 (ledger B18); (W4) and the field-dependent domains $\Lambda_j$ (ledger B20); (FX), (FX2) for the final step (ledger B19); (U) is the constant choice C23.

\paragraph{\file{SOURCES_D3R.md}: revision of D3 ((P3), (P8), freezing through R, Gevrey order).}
The sourced reset coefficient differs from Ba{\l}aban's by at most $\delta_r=4r+2C_hM\rho_1\lambda_k$, additively and pointwise (Lemma P3), so his sup bound of CMP 122 II p.~366 holds with one changed constant, and the large-field factors survive with $\gamma_0\to\gamma_0(1-\theta)\ge\gamma_0/2$, $\theta=\gamma^2\delta_r\le5\gamma^2r\le\tfrac12$. The t-parameter term splits as $d_B+d_S$ with $|d_S|\le2r$, freezing through R splits into a pure part (Lemma F1) and an R-induced part (Lemma F2), and the cumulant bound has Gevrey order $(n!)^{2+\epsilon}$, admissible with $D=3$.
Positivity is kept by a split against Ba{\l}aban's local coupling, and the (P8) term is split into a piece with the support of Ba{\l}aban's deviation, booked as a boundary term, and the source weight of D2. The regular output is split by an auxiliary parameter into the pure and R-induced parts, and the Gevrey count uses a Lah-number sum over partitions.
Verifier: \file{SOURCES_D3R_VERIFY.md}, \status{verified reduction} with corrections: (P3) needs $\rho_1\le r/(4C_hM)$; in (P8) $h_k$ is booked as irrelevant, and $d_S$ changes the covariance $C^{(k)}$ at every $t$ (relative size $O(\gamma^2r)$); Lemma F2 is a new claim, not Ba{\l}aban's bound for merged R-terms (item W8, taken up in \file{SOURCES_D3W8.md}); the Gevrey count is confirmed conditional on (B-N1).
Open inputs: strip holomorphy of $u_k$ (N3, ledger B18); (H-thr) (ledger B15); the real uniform (H-hist) (ledger B13); (W4) and the covariance change in (P8) (ledger B20, the latter by analogy with A16); (B-N1) (ledger B21).

\paragraph{\file{SOURCES_D3W8.md}: covariance change along $t$ and merged R-outputs (Lemma F2$'$).}
Along the t-interpolation the measure is Ba{\l}aban's CMP 119 (3.27) with covariance $C_f^{(k)}$, independent of $s$ and $t$; with $\tau\le3c_1c_H\gamma^2r\le\tfrac12$ one has $\|C_f-C_0\|_\delta\le2c_1\tau$, and the large-field factors survive with $\gamma_0\to\gamma_0/4$. Lemma F2$'$: the R-induced part obeys $a_m\le32\varepsilon'g_m^{\kappa_0}$ with no compounding, and $|E_R^{(m)}(X,U)-E_R^{(m)}(X,1)|\le CM^2\cdot32\varepsilon'(L^m\eta)^4g_m^{\kappa_0}e^{-\kappa'd_m(X)}$ on regular configurations, a marginal term of Ba{\l}aban's R-type form, summable in $m$ for $\kappa_0>2$.
Transported inputs keep their scale label and never pass through the generated map $G$, so a Schwarz-lemma step on $G-G(0)$ with Ba{\l}aban's factor $\varepsilon'=O(\alpha_2^{-1}\varepsilon_1)$ (CMP 109 p.~280) gives $a_{k+1}\le8\varepsilon'b_k$ per generation. The covariance bounds come from the weighted Neumann series of \file{SOURCES_D2.md}, and positivity from domination by a rescaled $C_0$ Gaussian.
Verifier: \file{SOURCES_D3W8_VERIFY.md}, route (a) and Lemma F2$'$ \status{proved+verified, conditional on (W4), (Eps), Lemma F1, (D), and $32\varepsilon'\le1$}, the t-measure and its bounds corrected: the cross term $g_k^2\langle A,\mathcal W_SA_k\rangle$ goes to $B^{(k+1)}$, CMP 116 p.~3 covers $C_f$ only by analogy, the domination must use the marginal region (or Anderson's inequality) to avoid a volume factor, and $\varepsilon'$ contains Ba{\l}aban's large constant of p.~280.
Open inputs: (W4) and (D) (ledger B20; (W4) is reduced to the printed (C1)--(C5) through Theorem SM); (Eps), from Theorem SM and A10; the extension A16 to the fluctuation covariance; the conditions $32\varepsilon'\le1$, $S\ge1$, $\kappa_0>2$ change the construction (ledger C16); the monotone-history assumption of Lemma F2$'$ is removed by the running maximum, with side condition $4\gamma^{\kappa_0}\le E_0$ (\file{SOURCES_GGAIN_NONMONO_VERIFY.md}).

\paragraph{\file{SOURCES_D4D5.md}: RP and symmetry transfer, and the cumulant bound (S$_{\mathrm{src}}$).}
The sourced RG output equals the bare integral $Z_\varepsilon[f]$ (Lemma D4.1), reflection positivity and the lattice symmetries are properties of the lattice measure (Lemmas D4.2--D4.4), and given (FX) the cumulants of real sources obey (S$_{\mathrm{src}}$): $|\kappa_n(f_1,\dots,f_n)|\le(C_S\sigma^{-a})^n(n!)^{1+\eta}\prod_iN(f_i)$ with $N(f)=\max(\|f\|_\infty/r,\|\nabla f\|_\infty/\rho_1)$, for $\sigma\in[16\varepsilon,\ell/8]$, supports pairwise at distance $\ge\sigma$, minimal diameter $\le\sigma$, and $C_S$ independent of $K$. Then (S$_{\mathrm{src}}$) $\Rightarrow$ (H-U) with $s=6$, $D>1$.
Each RG operation is an identity of integrals (Haar invariance, $\delta$-function blocking, partitions of unity, convergent expansions, $\int dV\,R\rho=\int dV\rho$), and complex sources follow by the identity theorem. Lemma C takes one Cauchy derivative in each of $n$ separated source amplitudes on a polydisc whose radius does not depend on $n$, so the sum over polymer sizes gives a single $n!$.
Verifier: \file{SOURCES_D4D5_VERIFY.md}, D4.1, D4.3, D4.4 and Lemma C \status{proved+verified} (Lemma C given (FX)), (S$_{\mathrm{src}}$) $\Rightarrow$ (H-U) confirmed, and (FX) confirmed as a new hypothesis not contained in CMP 116--122; corrected: cite CMP 122 I (1.102) for the actual R$'$ operation, the block geometry is invariant only under translations by $M\mathbb Z^4$ or coarser, and the order $1+\eta$ needs the coarsest-scale (FX2) radius (else an extra factor $(n!)^{2p_R/p_0}$, still OS-admissible).
Open inputs: (FX) in Ba{\l}aban's main form with the coarsest-scale (FX2) radius (ledger B19; remaining input N3, ledger B18); S1 at all scales, part of (B-S) (ledger B11); the (E14-c) clause (ledger B21).

\paragraph{\file{SOURCES_HBETA.md}: Lemma H, the two-loop bound (H$\beta$).}
Lemma H: given (I), (E), and (H-hist), with $x:=C_1A\gamma/\beta\le\ln(3/2)$ and $3r<g^{-2}-\gamma^{-2}$ ($g=g_K$), every real $c$ with $|c|\le2r$ keeps all $g_k(c)\in\ioc{\gamma}$ and gives $\sum_{j\le K}|B_j'(c)|\le e^x-1\le\tfrac12$ and $\tfrac12\le|v_k(c)|\le\tfrac32$, uniformly in $K$. With $\beta=\beta_0/2$ the smallness condition reads $\gamma\le\ln(3/2)\,\beta_0/(2C_1A)\approx0.20\,\beta_0/(C_1A)$.
In $u=1/g^2$ the chain rule gives $\partial\beta_j/\partial u_i=-\tfrac12g_i^3\,\partial\beta_j/\partial g_i$, so a $C^1$ bound costs $O(g_i^3)$ per step, $\sum_ig_i^3\le2\gamma/\beta$ along the trajectory, and a Gronwall product on the variational equation gives $e^x-1$. A bootstrap on the component of $\{|c|\le2r:u_K(c)>\gamma^{-2}\}$ that contains $0$ keeps every $c$-trajectory in $g\le\gamma$; the Gronwall algebra is formalized in Lean (\texttt{LemmaH}).
Verifier: \file{SOURCES_HBETA_VERIFY.md}, \status{proved+verified, conditional on (I), (H-hist)}; corrected: the hypothesis list adds that $\beta_j$ is the infinite-volume function independent of $K$, (H-hist) can be weakened to a column-sum bound, the smooth-clamp remark is dropped, and the toy numbers are not reproducible (the bound holds and is nearly sharp).
Open inputs: (I) (ledger B12; it holds by construction with the adapted gauge-fixing weight C10 and the frozen thresholds C1, but is not written as a lemma); the real uniform (H-hist) (ledger B13, a verified reduction with open inputs (Q)/$c_-$, ($M_r$), CMP 102 (189), and the (M-B) bookkeeping); (E), the uniform positivity of the one-loop $\beta_j(0)$; complex $c$ is not covered and is the strip bound N3 (ledger B18).

\paragraph{\file{SOURCES_T2_KAPPA2.md}: ledger B22.}
The decay loss $\kappa\to\kappa/2$ of the sourced step occurs at every step in three places (the interpolated inputs of the freezing proof in \file{SOURCES_D2.md} \S6.3, the extraction remainders of \file{SOURCES_D2C.md} \S3, the R-induced terms of Lemma F2 in \file{SOURCES_D3R.md}), and it is recovered by $L/2$-rescaling without compounding. The conditions are (K1) $\delta\le\tfrac1{10}(1-2(1+4\beta)/(\theta L))$, which is $\delta\le0.0385$ at $L=13$, $\beta=\tfrac14$, $\theta=\tfrac12$, and (K2) $\theta(1-\beta)L/2\ge1+3\beta$; at the last scale the final expansion runs at $\kappa_{\mathrm{eff}}=\theta\kappa/(1+4\beta)=\kappa/4$, with the $M$-condition (K3) $M\ge12L(4.5\kappa/C_1)^{1/4}$, implied by the existing one of \file{SOURCES_FX_R2.md}.
Ba{\l}aban proves only $2d_k(Z)\ge Ld_{k+1}(Z')$ (CMP 116 (2.36)), so the gain per step is $L/2=6.5$, and his printed choice $(1-10\delta)\tfrac12L=1$ in (2.41) spends all of it; the recovery therefore takes $\delta$ smaller, at the cost of a $\kappa$ about $3.6$ times larger, which moves only constants fixed after $\kappa$. The R operation drops $4\beta\kappa$ of rate (CMP 122 II (1.66)--(1.100)), the rescaling pays these drops and the T-step fraction $10\delta$, and the (1.93) volume charge is paid by $p_0$ through the condition before (1.95).
Verifier: \file{SOURCES_T2_KAPPA2_VERIFY.md}, \status{verified reduction}; items 2--4 confirmed on page images, item 1 corrected: the three losses stack within one step, so a literal half per absorption would make (K1) fail at $L=13$, and the repair reads $\theta:=1-m_{\mathrm{tot}}/\kappa$ as the kept fraction per cycle (each loss is a polynomial absorption with additive cost; $\theta\ge\tfrac12$ for $\kappa\ge2m_{\mathrm{tot}}$), with which (K1)--(K3) stand and the new lower bound is $\kappa\ge1536$ at $d=4$, $\beta=\tfrac14$.
Open inputs: the unprinted large-field T-step rate loss of CMP 119 Sect.\ 3 (ledger A23), on which the homogeneity of the large-field T step in the input rate rests; CMP 116 (2.36), which Ba{\l}aban asserts without proof.

\paragraph{\file{SOURCES_UN1.md}: (U), ledger C23.}
The scale condition (U) does not conflict with the bridge (B-N1): with $\lambda_K=1$ (physical units), the $r$ in (U) is the Cauchy margin $r_0$ of the constant-source family, not a test-function amplitude, so (U) is the constant choice $\rho_1\le c_*r_0/M^a$, made once with $\rho_1:=\min(c_*r_0/M^a,\,r_0/(4C_hM))$, after $r_0$ and before $K$. The bridge evaluates the generating functional only on sources rescaled into the domain of (FX2), at the price of a factor $(4/\rho_1)^n$ per $n$-point bound, with $s$ and $D$ unchanged.
The bridge uses $n$-linearity, $\kappa_n(f_1,\dots,f_n)=\prod_i\lambda_i\,\kappa_n(f_1/\lambda_1,\dots,f_n/\lambda_n)$, with $\lambda_i=2N_U(f_i)$ and $N_U(f)=\max(\|f\|_\infty/r,\|\nabla f\|_\infty/\rho_1)$, and pairwise $\sigma$-separated supports make the sup and the gradient of $\sum_iz_if_i$ a maximum over single terms, so the whole Cauchy polydisc stays in the admissible domain. Small widths $h\ge\varepsilon$ are paid by the norm growth $1/h$ per bump, which the flatness of $F$ absorbs, while (U) costs only an $h$-independent factor; the scale-matched reading of \file{SOURCES_D2.md} l.\ 217 would create a real conflict and is dropped.
Verifier: \file{SOURCES_UN1_VERIFY.md}, \status{proved+verified, conditional on the (H$\beta$) gradient bound}, so (U) is a constant choice (ledger C23), verified compatible with (B-N1); corrections: the (FX2) gradient radius is $\rho_1/2$, hence $\lambda_i=2N_U(f_i)$, the sup radius of the activity polydisc is $\tau_m(\Gamma)\le r$, overlapping supports never reach (S), and the (H$\beta$) bound $|\nabla f_k|\le\tfrac32\rho_1$ is a verified reduction for real $c$ (Lemma H and the chain rule), not open.
Open inputs: (H$\beta$) for complex $c$ (open); for real $c$ its inputs (I) (ledger B12) and (H-hist) (ledger B13); the $\kappa\to\kappa/2$ recovery (ledger B22, a verified reduction that uses A23).

\subsection{Step 3: the (S) clauses and (FX) (ledgers B21, B19)}\label{sec:src-FX}

\paragraph{\file{SOURCES_B21_S.md}: uniformity of (S), ledger B21.}
The source bound (S), $|\kappa_n(f_1,\dots,f_n)|\le(C_S\sigma^{-a})^n(n!)^{1+\eta}\prod_iN_U(f_i)$ for real sources with pairwise separation $\sigma\ge16\varepsilon$ and $\min_i\operatorname{diam}\operatorname{supp}f_i\le\sigma$, holds uniformly in the bump width $h\in[\varepsilon,\ell/512]$ and in the position of the bumps, with the C$^1$ norm $N_U(f)=\max(\|f\|_\infty/r,\|\nabla f\|_\infty/\rho_1)$. Every factor of $C^n(n!)^{1+\eta}$ is independent of $h$. The file traces one rescaled bump through Lemma C of \file{SOURCES_D4D5.md} and finds one $h$-dependent term, the cube count $Q_m\le(2+D/(M\lambda_m))^4$ at scales finer than the bump. With $D\le\sigma$ this count is at most $(2+2D/\sigma)^4u^4$, $u=\sigma/(2M\lambda_m)$, and the tree decay $e^{-3\kappa u/8}$ of the polymers that meet all supports cancels it, while the gradient cost is paid by $N_U$ and the flatness of the test function. Verifier: \file{SOURCES_B21_S_VERIFY.md}, \status{verified reduction}, corrected in scope (the inputs must also list the residue of B19 and D2 \S4.3) and in one constant (in the bridge $D/\sigma=3/8$, not $1/16$, so the factor is $(2.75)^4\approx57$). Open inputs: S1 at all scales including $m=0$; (FX) in Ba{\l}aban's main form (1.98) (ledger B19 with the choice C13), whose remaining input is N3 (B18); the coarsest-scale (FX2) radius, without which the order is $2+\eta$ and $D>2$; D2 \S4.3, itself only reduced; the (E14-c) clause of B21, which no file maps to Lemma C.

\paragraph{\file{SOURCES_FX.md}: hypothesis (FX), ledger B19.}
(FX), the convergence of the final Mayer expansion of the sum over large-field geometries plus the unit-scale integral, with polymer radius $c\cdot r$ in the source (FX2), uniform in $K$ and $\varepsilon$, is reduced to Ba{\l}aban's large-field bounds CMP 122 II (1.79)--(1.89) and (1.97), the convergence criterion of CMP 116 (2.38)--(2.41) and Cammarota (ledger X4), and named local checks. The construction is a final R-step with $Z''=\emptyset$ followed by the unit-scale cluster expansion of CMP 116 \S2, as Dimock does for $\varphi^4_3$ (Dimock III, Theorem 2). The file reads CMP 122 I (0.2)--(0.6) and CMP 122 II (1.71)--(1.104) and finds that Ba{\l}aban already exponentiates the sum over renormalizable large-field components inside each R step, so that only the components that pay for later steps stay as a sum, and at the last scale there are no later steps. The fine-history sum inside a component is bounded by (1.79), because Dimock's power-decay argument ((209)--(211)) does not transfer to $d=4$, and the source enters the large-field operations only through the upper bound (1.73). Verifier: \file{SOURCES_FX_VERIFY.md}, \status{verified reduction}, corrected: the per-step expansion sums only components of diameter $\le100MR_k$ counted by position ((1.96)), so the last step, which sums components of unbounded size, needs the tree decay O1; three page quotes are fixed ((1.80) has $R_n^{d+1}$ and a local $K$, (1.85) has $A_1^2p_0^2$, the condition after (1.94) ends in $-\tfrac32p_0(g_k)$); the fallback order is at most $1+\eta+2/3$. Open inputs: the verifier left O1(b), O2, and O3 open; O1(b) and O2 are now \status{proved+verified} (\file{SOURCES_FX_O1.md}, \file{SOURCES_FX_O23.md}), R2 is \status{proved+verified} as a reduction, and O4 is avoided by the main form (C13), so the remaining input of B19 is N3 (ledger B18), through O3.

\paragraph{\file{SOURCES_FX_O1.md}: gap O1(b) of (FX), Lemma [P].}
Lemma [P]: for a connected union $Y$ of $MR_j$-cubes, $d'_{j+1}(S(Y))\ge c_*d'_j(Y)$ with $c_*=[3\cdot4^d\sqrt d(\rho+2)^d]^{-1}$, $\rho=LR_{j+1}/R_j$, and no additive constant. With the $n=K+1$ term of CMP 122 II (1.80) this gives $\kappa_K(Y)\ge\kappa_{\mathrm{fin}}d_K(Y)$ for every non-renormalizable final component $Y$, which is Ba{\l}aban's assertion of the term $-\kappa_1d_k(X)$ on p.\ 387, with $\kappa_1=\kappa_{\mathrm{fin}}$. The proof combines three elementary facts (a tree of length $\ell$ meets at most $4^d(1+2\ell)$ closed unit cubes, $d'_j(Y)\le\sqrt d(n_f-1)$, and a coarse cube meets at most $(\rho+2)^d$ fine cubes) with the count of at least $21^d$ coarse cubes in $S(Y)$. The halo $S(Y)\setminus Y$ carries no large fields (p.\ 384), so nothing is absorbed by $p_0$, and only the first term of (1.80) is used, all its terms being non-negative. Verifier: \file{SOURCES_FX_O1_VERIFY.md}, \status{proved+verified}, with one remark corrected (the counterexample to the form ``factor $1/L$ plus a constant'' needs $d\ge2$, a constant independent of $L$, and $L$ large) and R3 discharged by Ba{\l}aban's own comparison in (1.86)$\to$(1.87). Open inputs: R2 ($\kappa_{\mathrm{fin}}\ge\kappa_{\mathrm{final}}$), the only point left by the verifier, is closed in \file{SOURCES_FX_R2.md} (choice C14); R1 (grid alignment) changes constants only.

\paragraph{\file{SOURCES_FX_O23.md}: gaps O2 and O3 of (FX).}
O2: the R-denominator of CMP 122 I (1.100) is one small-field integral per renormalizable component with the bare Wilson action $-g_k^{-2}A(\zeta_i,\cdot)$ and no E-, R-, or B-terms, so it is holomorphic and non-vanishing on a source polydisc of radius $c\cdot r$, with $c$ independent of $|\Lambda_i|$, $k$, and $\varepsilon$. O3: inside the large-field operations $T'$ the sourced weight obeys $\operatorname{Re}(w+\delta w)\ge(1-\theta_W)(g''_k)^{-2}$, so the upper bounds (1.73), (1.79), (1.89), (1.92), (1.97), (1.100) hold with $\gamma_0\to\gamma_0(1-\theta_W)(1-\tau_C)\ge\gamma_0/4$ and $p_0\ge16/(\gamma_0A_1^2)$. For O2 the Wilson part of the source perturbation is $|\sigma_W|\le\tau\Phi(g_k)$ with $\Phi(g)=C(M)g^2p_0(g)^2R(g)^4\to0$, and (1.102) also allows an $f$-independent denominator. For O3 the plaquette terms $P_p\ge0$ of the real $G$-valued configuration inside $T'$ (CMP 122 II p.\ 379) give domination in modulus by a positive operation with a smaller coupling prefactor, while the thresholds stay fixed. Verifier: \file{SOURCES_FX_O23_VERIFY.md}, \status{proved+verified} for O2 and \status{proved+verified, conditional on N3} for O3, since the strip bound $|\delta w|\le\tfrac32\tau$ is N3; it adds the cleaner condition $p_0\ge8/(\gamma_0A_1^2)$ with $\gamma_0'\ge\tfrac34\gamma_0$, and corrects O4 (the division by $T'_k(Y_i)1$ in (1.103) has no positivity proof, but (1.103)--(1.104) is an optional final remark, so O4 is avoided by keeping the main form (1.98)). Open inputs: N3 (ledger B18); the main form is the choice C13; O4 binds only if (1.104) were used, with fallback order $2+\eta$.

\paragraph{\file{SOURCES_FX_R2.md}: gap R2 of (FX), $\kappa_{\mathrm{fin}}\ge\kappa_{\mathrm{final}}$.}
R2 holds as a reduction in volume form: $\kappa_K(Y)\ge\kappa_{\mathrm{vol}}\#M(Y)$ with $\kappa_{\mathrm{vol}}=\sqrt d\,\kappa_{\mathrm{fin}}=C_1M^dR_{K+1}/(3\cdot4^d(L+2r)^d)$, $r=R_K/R_{K+1}$, and $\kappa_{\mathrm{vol}}\ge\kappa_{\mathrm{final}}:=2\kappa+\sqrt d(1+2\beta)\kappa+c_d$ once $M\ge12L(3\kappa_{\mathrm{final}}/C_1)^{1/d}$, uniformly in $K$, $R$, and $g$. Here $C_1:=\max(1,C_c)$ is the enlarged constant of (1.80). The file quotes Ba{\l}aban's order of constants ($L$ first, $\kappa\ge\kappa_0$ before $M\ge M(\kappa)$, $p_0$ and $\gamma$ last), so $\kappa_{\mathrm{final}}=\kappa_{\mathrm{final}}(d,\kappa,\beta)$ is fixed before $M$ (route (a)), after $C_1$ is enlarged (route (b)). It replaces the tree rate of \file{SOURCES_FX_O1.md} by the volume rate, because the final expansion charges large-field domains per $M$-cube and the tree rate alone loses $R_K^{d-1}$. Verifier: \file{SOURCES_FX_R2_VERIFY.md}, \status{proved+verified} as a reduction, corrected in its reasoning: the $2\kappa$ is the (1.93) decay-recovery charge $(1+\beta+\beta/(3\cdot2^d))\kappa$ per $M$-cube, not a Mayer-link cost, and the term $\sqrt d(1+2\beta)\kappa$ is surplus. Open inputs: none; the two minor inputs, $R_{K+1}\ge1$ and $C_c$ independent of $g$, are verified readings (former ledger A20, \file{SOURCES_BOOKKEEPING.md}), and $\kappa_{\mathrm{final}}$ is the choice C14, not sharp.

\subsection{Steps 4--5: N3, (H-hol), and the large-field inputs (ledgers B18, B16, B17)}\label{sec:src-N3}

\paragraph{\file{SOURCES_HBETA_N3.md}: N3, the strip bound, ledger B18.}
With the source additive in $u=g^{-2}$ ($w_0=g_0^{-2}-f$, $\tau$ in absolute $u$-units), N3 is equivalent to Lemma H of \file{SOURCES_HBETA.md} for complex $c$ on $Q=\{|\operatorname{Re}c|\le2r,\ |\operatorname{Im}c|\le\tau\}$, pointwise in $x$. Lemma N3 gives $|\delta w_k|\le\tfrac32\tau$ for all $k\le K$, uniformly in $K$ and $\varepsilon$, given (I), (E), real (H-hist), (H-hol) on relative polydiscs, and (B-hol), with (H2-hist) added in the history form; the constant $\tfrac32$ can be any $1+\eta$. Step 0 bounds $\partial\beta_j/\partial u$ on the thin strip by Ba{\l}aban's real C$^1$ bound at $\operatorname{Re}u$ plus a Cauchy correction of order $g^4|\operatorname{Im}u|$, and a complex Gronwall step with $\sum_i(\operatorname{Re}u_i)^{-3/2}\le4\gamma/\beta$ closes a bootstrap in $\operatorname{Im}c$. The file also shows that Cauchy estimates on relative discs alone give a factor $1+Cg^2$ per step and a bound that grows with $K$, so the real derivative bound is needed. Verifier: \file{SOURCES_HBETA_N3_VERIFY.md}, \status{verified reduction}, corrected: condition (a), $3r+\tfrac32\tau<g^{-2}-\gamma^{-2}$, limits $r$ by the gap between $g=g_K$ and $\gamma$ and does not follow from $\gamma$ small (settled in \file{SOURCES_N3A.md}); $d\beta/du=O(g^3)$, or $O(g^4)$ for even $\beta$; (H2-hist) is sufficient, not necessary, the minimal input being (H-hist$_C$) on the complex polydisc. Open inputs: (I) (ledger B12), (E), real (H-hist) (B13), (H-hol) in its small-field (B16) and large-field (B17) halves, (H2-hist) or its Route 2(a) (B14), (B-hol), the $h\equiv0$ booking (C11), and the additive normalization (C12); Step 3 in the D$''$ form is proved, unverified (B10).

\paragraph{\file{SOURCES_N3A.md}: condition (a) of Lemma N3.}
With the order of constants $L\to$ Ba{\l}aban's constants $\to\gamma\to g\to(r_0,\tau_0)\to K$ and $g=g_K\le\min(\gamma/2,g_1)$, condition (a) holds at the uniform radius $r_0=\tau_0=\min(1/8,r_0^{\mathrm{D2}},r_0^{\mathrm{D3}},r_0^{\mathrm{FX}},\gamma^{-2}/10,1/(6\gamma^2c_1c_H))$, since $\tfrac92r_0\le\tfrac9{16}<3\le g^{-2}-\gamma^{-2}$ for $\gamma\le1$, a margin of $5.33$. The radius does not depend on $K$, and in $u$-units it is the physical sup-norm radius of the source. The file checks that $g\le\gamma/2$ is admissible for every $K$ (\file{LEMMA_DPRIME.md} Step 8 in the Lemma E form, \file{KAPPA_BOUNDED.md} \S5.3(b) in the D$''$ form) and that nothing forces $g_K$ near $\gamma$. It refutes the proposed radius $r=\tau=\gamma^{-2}/2$, which satisfies (a) but breaks conditions (b)--(d) of Lemma N3 and the radius bounds of D3R and D3W8. Verifier: \file{SOURCES_N3A_VERIFY.md}, \status{proved+verified}, with item 1 corrected: ``$u_k$ decreasing, $g_K$ largest'' holds only in the Lemma E form, while in the D$''$ form only $u_k\ge u_K-D'$ holds and $g_0$ can exceed $g_K$; the D$''$ sketch has two unnamed gaps. Open inputs: none for condition (a); the order of constants is the choice C5; the inputs of N3 are unchanged, and N3 in the D$''$ form stays a sketch (ledger B10).

\paragraph{\file{SOURCES_HHOL.md}: (H-hol), small field, ledger B16.}
For the small-field chain, $\beta_j$ extends holomorphically to the polydisc $|g_i-x_i|<\rho_0x_i$ with $|\beta_j|\le M$ uniformly in $j$, $k$, and volume, where $\rho_0=\min(0.1,(LM)^{-2}(8n)^{-1/2})$, divided by $\sqrt2$ when the $X$-modes are counted (ledger C6). This holds given Ba{\l}aban's conditions with modified constants: (C1$'$)--(C4$'$) with $\varepsilon_1\to1.235c\varepsilon_1$, (C4$''$) $8n\rho_0^2(LM)^4\le1$ (with $16n\rho_0^2$ once the $X$-modes are counted), and (2.24$'$); the radius $\rho\le1/3$ first asked for is not reached, and Lemma N3 accepts any fixed $\rho_0$. In reference units $\tilde B=xB$ the complex coupling enters only through $w=(x/g)^2$, as the scale of the Gaussian and as the coefficient of the $u$-prefactored vertices, while the cut $|\tilde B|<\varepsilon_1$ does not depend on $g$. The marginal-domination lemma of the file (its Lemma D) bounds the complex Gaussian by the real one at $\operatorname{Re}w$ at a cost $(|w|/\operatorname{Re}w)^{1/2}\le e^{2\rho^2}$ per real mode, so the volume-uniform bound is made per polymer inside CMP 116 (2.14)--(2.17). Verifier: \file{SOURCES_HHOL_VERIFY.md}, \status{verified reduction}, corrected: the per-polymer cost misses the $X$-modes ($\rho_0\to\rho_0/\sqrt2$), $|\beta_j|\le M$ is a reduction because the majorant covers the Gaussian only, and (C1)--(C5) are not Ba{\l}aban's printed labels; single-valuedness is confirmed with no extra assumption. Open inputs: the modified conditions (C1$'$)--(C4$'$), (C4$''$), (2.24$'$) (the $X$-mode question is resolved in the row M-C of the dependency table); (F1), reduced to Theorem SM at complex couplings, from which (H2-hist) is derived with $a_m=\lambda_C^{m/2}$; a $\gamma$-independent $C_1$ in real (H-hist) (B13); the large-field half (ledger B17, \file{SOURCES_HHOL_LF.md}).

\paragraph{\file{SOURCES_HHOL_LF.md}: (H-hol), large field, ledger B17.}
The large-field half of (H-hol) holds on $P_{\rho_0}(x^*)$ around Ba{\l}aban's real trajectory $x^*$, with the R-ratio density, the thresholds, $R_k$, and $N$ frozen at $x^*$ (ledger C1): the large-field integrals are entire in $u=g^{-2}$ in finite volume and are dominated in modulus with $\gamma_0\to\gamma_0''\ge0.82\gamma_0$ at $\rho_0=0.1$, and F2$'$ holds on the polydisc with $32\varepsilon'\to40\varepsilon'$. This is a reduction to (M-A), (M-B), (M-C), and (M-D). The file shows that the radius $(100MR_k)^{-4}$ of W4 is a radius in the input activities, which the real (1.100) denominator does not contain, and it removes the coupling from the denominator by freezing the ratio at $x^*$, at a mismatch cost of $O(1)$ per unit cube inside Ba{\l}aban's budget $p_0(g)\gtrsim R(g)^{d+2}$. F2$'$ compares the bound sizes $(x^*_j)^{\kappa_0}$, which are ordered on the real trajectory, so no ordering of complex couplings is needed, and a slot-dependent radius $\rho_j\propto x_j^2$ would have failed Lemma N3. Verifier: \file{SOURCES_HHOL_LF_VERIFY.md}, \status{verified reduction}, corrected: $R_k$ is unbounded, of order $(\log K)^r$, but largest at $k=0$ and non-increasing in $k$ (the file reverses the index); $\theta=0.174$, not $0.26$; (M-D) is not a convention, since the frozen map equals Ba{\l}aban's map only at $x^*$ and its derivative there is (H-hist) plus (H-thr). Open inputs: (M-A), reduced to (H$-$) and CMP 102 (189) (ledger A4, with A18 and A19), with (M-A$'$) closed; (M-C), a verified reduction with open (I-$\varepsilon$), the (I-V0) residual, and (M1); (M-D), which needs real (H-hist) (B13) and (H-thr) (B15) with (H-hist2) (B14); (M-B) is closed, \status{proved+verified} with $p_0\ge17/(\gamma_0A_1^2)$ and $\theta_W+\tau_C\le0.2$.

\paragraph{\file{SOURCES_MA.md}: (M-A), the equipartition bound.}
For the cut-off ratio density of CMP 122 I (1.100) at real weight $0<u\le x^{-2}$, the file proves $\langle A-A_{\min}\rangle_u\le(n_i/u)(1+\eta_x)$, where $n_i=(N^2-1)L_i$ counts the real coordinates on the free links and $\eta_x=0.68M_0^2x^2p_0(x)^2+2c_HM_0^2x\,p_0(x)^3$; the same bound with the full $A$ is false in general, because the background minimum $A_{\min}$ is $O(\varepsilon^2)$ per block. As a corollary, the freeze-mismatch average is at most $e^{0.1906\,n_i(1+\eta_x)}$, that is $e^{0.77(N^2-1)(1+\eta_x)}$ per scale-$k$ unit cube ($e^{2.3}$ for SU(2)), additive in the volume.
A virial identity on the convex cutoff ball, centred at the true minimizer $X^*$ of the exponent, bounds $u\langle(X-X^*)\cdot\nabla G\rangle_u$, $G=A-A_{\min}$, by $n_i$ plus a Haar-Jacobian term (Lemma J: $0.68N(N-1)r^2$ per link for $r\le\tfrac12$). A segment bound from the Hessian lower bound (H$-$) (Lemma S) turns this into the bound on $\langle G\rangle_u$, and integrating $\tfrac{d}{du}\log D^G=-\langle G\rangle_u$ gives the corollary.
Verifier: \file{SOURCES_MA_VERIFY.md}, \status{proved+verified, conditional on (H$-$)}; corrections: the proved constant is $C=1$, not $\tfrac12$ (sharp when $X^*$ lies on the boundary), the SU(2) toy value 1.87 needs a cutoff norm that the file does not state, and the bookkeeping item (M-A$'$) is a real gap, closed in \file{SOURCES_MA2.md}.
Open inputs: (H$-$), the Hessian lower bound along segments with $c_H$ uniform in $k$, volume, and background; \file{SOURCES_MA2.md} reduces it to CMP 102 (189) (ledger A4) together with A18, A19, and its conditions (C1), (C2); ledger B17.

\paragraph{\file{SOURCES_MA2.md}: (M-A$'$) and (H$-$).}
The bookkeeping item (M-A$'$) closes: every normalized ratio of CMP 122 I (1.100) sees only $G=F-A_{\min}$, and the exact identity $N_x/D_x=e^{-uF}\,[e^{uA_{\min}}/D^G(x^{-2})]\,e^{(u-x^{-2})G}$ keeps $e^{uA_{\min}}$ at the coupling $u$, so the mismatch costs at most $e^{1.07(N^2-1)(1+\eta_x)}$ per unit cube. The Hessian bound (H$-$) of \file{SOURCES_MA.md} is derived in the form $D^2F[h,h]\ge-c_H\varepsilon_k\sum_b\|h_b\|_2^2$, with $c_H$ uniform in $k$, volume, and regular backgrounds.
Since $A_{\min}$ does not depend on $V'$, $D^A(s)=e^{-sA_{\min}}D^G(s)$, and the error in \file{SOURCES_HHOL_LF.md} \S3.4(c) was a numerator at $u$ over a normalizer at $x$. For (H$-$), the chain rule per plaquette is combined with one Cauchy step on $\delta\mathcal H/\delta B$ minus its linear part on the complex ball of CMP 102 (172), where the $1/\varepsilon$ of the radius cancels the $\varepsilon$ of CMP 102 (187), and the kernel decay of (189) closes a Schur test.
Verifier: \file{SOURCES_MA2_VERIFY.md}, (M-A$'$) \status{proved+verified} and (H$-$) a \status{verified reduction} to CMP 102 (189); corrections: the first Hessian term is $\ge0$, $c_N=O(1)(B_0C_4+C_2)B_5$, $c_H$ depends on $M^5$ and $M_0$, and $A_{\min}$ may be only piecewise analytic in the background, so $e^{uA_{\min}}$ must stay paired with an analytic factor or inside modulus bounds.
Open inputs: (C1) $M_0\varepsilon_k\le C_1\varepsilon_1$; (C2) the localized map $U_{k,X_i}$ obeys the CMP 102 bounds; (C3) CMP 102 (189), announced without proof and already used in Ba{\l}aban's own chain (ledger A4, now a verified reduction to A1); ledger A18 (including the $k$-uniformity of $C_2$, $C_4$, $B_0$) and A19; and the reading that Ba{\l}aban's chain removes $e^{uA_{\min}}$ by algebra.

\paragraph{\file{SOURCES_MB.md}: (M-B), the large-field bookkeeping.}
Ba{\l}aban's large-field bookkeeping CMP 122 II (1.73)--(1.99) survives the complex-coupling perturbation with thresholds frozen at the real trajectory $x^*$: the per-cube factors $\exp(-p_0(x^*_j))$ and the per-component factors $\exp(-2p_0(x^*_j))$ keep the same $\varepsilon_j$, $\delta_j$, $\delta_j'$, $R_j$, and the only new condition is $p_0\ge17/(\gamma_0A_1^2)$ with Ba{\l}aban's output factor kept (option B). The cost constant becomes $C_c'=C_c+C_m+C_G+C_E$, a fixed $O(1)$ per site inside the $O(1)\log g_j^{-2}$ budget, and the p.~389 condition $p_0\ge8\kappa(100R_k)^d+4$ is unchanged.
Each line of (1.73)--(1.99) that carries $\gamma_0$, $p_0$, $\varepsilon_j$, $\delta_j$, $A_1$, or a per-cube factor is listed (L1--L24) and redone with $\gamma_0\to g^*\gamma_0$, which is legitimate because each factor bounds a positive Wilson action or quadratic form whose modulus at complex $u$ is dominated with prefactor $\mathrm{Re}\,u\ge0.8264\,(x^*)^{-2}$. The combined factor of p.~381 then needs $p_0\ge c/(\gamma_0A_1^2)$ with $c=2/(g^*-\tfrac12)$, and the budget constant $C_1'=\max(1,C_c')$ (route (b) of \file{SOURCES_FX_R2.md}) multiplies (1.82), (1.88) by a fixed factor that $p_0$ large and $\gamma$ small absorb.
Verifier: \file{SOURCES_MB_VERIFY.md}, \status{proved+verified, conditional on (M-A) and (M-C)}; corrections: the perturbation weights are additive, so $g^*=0.82645-\theta_W-\tau_C$ (not a product), and the constant 17 needs the margin $\theta_W+\tau_C\le0.2$ (strict limit 0.2088; at $\tfrac14$ one needs $c=26.2$); p.~383 prints $\tfrac12\gamma_0$ and $N^{\beta_0}$, and any fixed $g^*>0$ is accepted there.
Open inputs: (M-A) (the freeze-mismatch cost $C_m$) and (M-C) (the complex-$u$ activities $\alpha^c$ and E-terms); the link from F2$'$ to $\alpha$ stays an inference; the margin is ledger C2, and the block feeds ledger B17.

\paragraph{\file{SOURCES_MC.md}: (M-C), the complex small-field step.}
Lemma MC: if Ba{\l}aban's constants satisfy the conditions B1--B23 of CMP 116 \S2, then with the shrinks $\varepsilon_1'=\varepsilon_1/(4\lambda_V)$, $\gamma'=\gamma/(8\lambda_V)$, $\alpha'=\alpha/4$, $\lambda_V=\max(c,1.235)$, the radius $\rho_0=\min\{0.1,(LM)^{-2}(0.45a/(16n))^{1/2}\}$, and the added condition (M1), the small-field step is holomorphic in complex $g$ with $|g/x-1|<\rho_0$, for every real $x\le\gamma'$ and inputs in the ball $cE_0$, and Lemma 3 (2.38) holds with Ba{\l}aban's $C_3\varepsilon_1$. Hence $\beta_j$ is holomorphic on these discs and bounded uniformly in $j$, $k$, and volume.
The 23 inequalities of CMP 116 \S2 are read on the page images and typed, and each one at complex $w=(x/g)^2$ and primed constants is shown to follow from Ba{\l}aban's own. The vertex is affine in $w$ with real-domain coefficients, so its factor $\lambda_V$ is absorbed by the $\tau$-circle (2.18) with $E_0\to\lambda_VE_0$ and a smaller $\varepsilon_1$; the cut stays real, and the Gaussian cost $16n\rho^2$ per site and the $\alpha_5$ items fit into the budget $(LM)^4\alpha_i\le\tfrac14$ under (M1) and (M2).
Verifier: \file{SOURCES_MC_VERIFY.md}, \status{verified reduction} to the printed CMP 116 \S2 conditions (B1--B23 and four more), Lemma 2, (1.36), and p.~12 (2.5)--(2.6); corrections: the shrinks are $\varepsilon_1/(4\lambda_V)$ and $\gamma/(8\lambda_V)$, not $\varepsilon_1/\lambda_V$ and $\gamma/\lambda_V$, the total output is $\le E_0$, not $\tfrac12E_0$, and the open item (I-X) is resolved by p.~12 (the $X$-measure is the unit Gaussian).
Open inputs at the time of the verifier: (I-$\varepsilon$), that the smaller $\varepsilon_1$, $\alpha_0$, $\alpha_1$ keep CMP 116 \S1 and the Part I induction valid; the (I-V0) residual (the $\log\sigma$ and Jacobian part); and (M1), an ``$M$ large'' condition; ledger B17. All three are now settled, see the next four paragraphs: (I-V0), (I-ord), and (I-$\varepsilon$) in both layers are closed and verified, and (M1) is the convention C24.

\paragraph{\file{SOURCES_MC_IV0.md}: (I-V0), the vertex part of (M-C).}
The $w$-dependence of the bracket $\{\dots\}$ in CMP 116 pp.~9--10 is only the $E_k$ difference. It is linear in $w$ and free of $w$ otherwise. $\operatorname{Tr}\log$ and $\log\sigma$ sit in $P^{(k)}$ with prefactor 1, and $\log\sigma$ and the Jacobian do not depend on $w$. The five $1/g_k^2$ terms are each of the form $u\cdot f(C\tilde B)$, so they are affine in $w$, and the shift restores (2.19), (2.20), and (2.26) exactly.
The file reads pp.~9--11 and 16--17 on 300 dpi page images and checks each $w$-dependent term against the printed display.
Verifier: \file{SOURCES_MC_IV0_VERIFY.md}, \status{proved+verified} for the closure of (I-V0); corrections: the reading of the integral as a Haar integral is an inference [I], since the page prints only $dB'\sigma(B')$; the $\varepsilon_1^2$ cut factor of (2.26) belongs to (I-$\varepsilon$).
Open inputs: none for (I-V0); (I-lin) is a mild item (per-piece balls).

\paragraph{\file{SOURCES_T2_IEPS.md}: (I-$\varepsilon$) and (I-ord), small-field layer.}
Every condition on $\varepsilon_1$, $\alpha_0$, $\alpha_1$ in CMP 116 \S1--\S2 and CMP 109 \S1, \S3 is an upper bound, a ratio of degree one, or a lower bound relative to a later constant; none is two-sided. So the shrinks of Lemma MC keep these papers valid: the $\varepsilon_1^2$ cut of (2.26) is restored exactly by $\gamma'=\gamma/(8\lambda_V)$, and scaling $(A_0,A_1)$ by $1/\theta$ restores CMP 122 II pp.~380--381 exactly. (I-ord): the order $d,L\to\delta\to\kappa\to\kappa_1\to$ walk constants $\to M\to\beta,\alpha'\to\alpha_1\to\alpha_0\to\alpha_3\to\alpha_2\to E_0\to\varepsilon_1\to\rho_0\to A_1\to A_0,p_0\to\gamma$ has no cycle.
Verifier: \file{SOURCES_T2_IEPS_VERIFY.md}, \status{proved+verified} for CMP 109 and CMP 116 (five missed conditions found, all of an allowed type); \status{open} at that time for CMP 119/122, because CMP 119 (2.28) defines $\alpha_{i,j}=g_jC_i(\log g_j^{-2})^{q_i}$ with $C_0$, $C_1$ ``sufficiently large''; the restoration of (L1) by a smaller $\gamma$ in the large-field layer was unsupported. (I-ord) is closed on the pages read.
Open inputs: the large-field layer, closed in the next paragraph.

\paragraph{\file{SOURCES_T2_IEPS_LF.md}: (I-$\varepsilon$), large-field layer.}
The shrink of Lemma MC reaches $\alpha_{i,j}$ (CMP 119 p.~262 takes the restrictions from CMP 109). Every printed condition on $C_i$, $q_i$ is a lower bound, a ratio ($q_0\le q_1$), or an upper bound on $g$; none bounds $C_i$ from above. Lemma MC needs only a factor-4 slack, so $C_i$ and $q_i$ stay fixed and $\gamma$ is lowered. The order continues $\rho_0\to$ exponents $\to A_1\to A_0\to C_0\to C_1\to\gamma$, with no cycle. Residue [I]: $\alpha_2$--$\alpha_4$ and $\gamma_2$ are not printed in CMP 119/122 (all three branches close); $q_0>p_0$ versus $q_0=p_0$ is not printed (both close).
Verifier: \file{SOURCES_T2_IEPS_LF_VERIFY.md}, \status{proved+verified}: all four claims confirmed; corrections: the condition C3 is a printed inclusion with an unprinted constant $c(B_3)$, and $q_0\ge p_0$ is forced; the p.~381 factor is $\exp(-\gamma_0A_1^2p_1^2(g_j))$ with no extra weight; one inequality on $\nu$ reverses ($\tfrac12\nu\le p_0-p_1-1$); $\gamma'<\min(e^{-q_0},e^{-q_1})$ is a further upper bound on $\gamma$, harmless.
Open inputs: none beyond the two [I] residues above, each of which closes on every branch.

\paragraph{\file{SOURCES_HIGHER_M1.md}: higher derivatives, and (M1).}
Part A: no consumer in the coupling and threshold chain needs derivatives of order greater than 2 in the couplings or thresholds. Part B: the condition (M1), ``$M$ large'', is a block-size choice. The file lists 26 lower bounds on $M$ (for example $M\gg\kappa_1$, $M\ge M(\kappa)$, $M\ge cL$, $M\ge48\kappa_1/\varepsilon$, $M/M_1>3(6+c_X)$) and no upper bound, and the order of constants puts $M$ before $\alpha$. So (M1) is the convention C24: $M$ is chosen at least the largest of these lower bounds.
Verifier: \file{SOURCES_HIGHER_M1_VERIFY.md}, \status{proved+verified} with minor corrections: Part A is restricted to the coupling, threshold, and MA2-field variables (Ba{\l}aban's field expansion goes to order 4, and the source derivatives of D2/D3 to all orders); the lower bounds of \file{SOURCES_A3_SIGMA.md} 5.1 and 5.3 join the list.
Open inputs: none; (M1) leaves ledger B17 and becomes ledger C24.

\subsection{Step 6: thresholds and second history derivatives (ledgers B15, B14)}\label{sec:src-thr}

\paragraph{\file{SOURCES_HTHR.md}: (H-thr), first attempt.}
With the threshold mismatch $\Delta_j(c)=B_j(u(c);\bar u)-B_j(u(c);u(c))$ (integer parameters frozen, convention L3), the file bounds $\sup_c|\Delta_j'(c)|\le3(1+3r)A\,\Phi(g_j^{-2})$ and sums this along the trajectory to $C\beta^{-1}\gamma^{-2(1+a)}(\log\gamma^{-2})^P$ times a decay factor, uniformly in $K$; in the verifier's corrected form the decay is $\exp(-(\log\gamma^{-2})^q)$ with $q=2p_1-(d+5)r_0>p_0$, in place of $\exp(-\kappa(\log\gamma^{-2})^{2p_0})$. Lemma H$'$ then splits the budget: $x\le\ln(6/5)$ for (H-hist) and $x_{\mathrm{thr}}\le\tfrac14$ give $\sum_j|B_j'|\le\tfrac12$ and $\tfrac12\le|v_k|\le\tfrac32$.
Coarea writes $\Delta_j$ as the contribution of the shell between the two thresholds, which reduces $\Delta_j'$ to threshold derivatives and mixed threshold--weight derivatives of $B_j$, bounded in Lemma S under the model assumptions M1--M4. The sum along the trajectory uses $u_i\ge\gamma^{-2}+\beta m$ and an integral comparison that keeps the factor $\gamma^{-2}/\beta$, and Lemma H$'$ is an additive Gronwall argument.
Verifier: \file{SOURCES_HTHR_VERIFY.md}, \status{verified reduction} of (H-thr) to M2$'$, M4, M5, with the shell representation, the trajectory sum, and Lemma H$'$ confirmed; corrections to Lemma S: the fluctuation cutoff is $\delta_j'=g_jA_1p_1(g_j)$ with $p_1<p_0$, Ba{\l}aban bounds one bond only through the block quadratic form CMP 122 II (1.78) with loss $R_j^{-(d+5)}$, so the per-link density M2 is unsupported, and the mixed derivative needs history holomorphy (M5); M1 is resolved ($\chi$ stays inside the integral, so $\Delta_j\ne0$).
Open inputs: M2$'$ (density form of the block bound), M4 (extraction cost), M5 (history holomorphy), M3, and the frozen-integer convention L3 (ledger B15, C1); \file{SOURCES_HTHR2.md} replaces M2$'$ by Lemma T and M5 by (H-hist2).

\paragraph{\file{SOURCES_HTHR2.md}: (H-thr), Lemma T.}
The shell indicator is at most the large-field function $1-\chi$ at the lower threshold, so Ba{\l}aban's printed tail factor of CMP 122 II (1.77)--(1.78) (with the misprint $p_1\to p_1^2$) bounds $\Delta_j$ itself with exponent at least $\kappa_2(\log s)^q$, $\kappa_2=\tfrac12L^{-(d+5)}$, $q=2p_1-(d+5)r_0>p_0$ (Lemma T). With a uniform bound $D$ on $\Delta_j''$ this gives $\sup_c|\Delta_j'|\le2(2AD)^{1/2}\Phi_T(u_{j-1})^{1/2}+(2A/r)\Phi_T(u_{j-1})$, a sum over $j$ that is uniform in $K$ under (G1$'$)--(G3$'$), and Lemma H$'$ of \file{SOURCES_HTHR.md} unchanged.
The zero-order bound on $\Delta_j$ uses Lemma T, the history weights M3, and the extraction cost M4, a Cauchy estimate on the analyticity radius $\alpha_{1,j}=g_jC_1(\log g_j^{-2})^{q_1}$ of CMP 119 (2.28), so no history derivative enters. The derivative then comes from the interpolation inequality $\|\Delta_j'\|\le2(\|\Delta_j\|\,\|\Delta_j''\|)^{1/2}+\|\Delta_j\|/r$ (Lemma 7.1(a) of \file{SOURCES_D2.md}), which replaces the history holomorphy M5 by a uniform second-derivative bound (H-hist2).
Verifier: \file{SOURCES_HTHR2_VERIFY.md}, \status{proved+verified, conditional on (H-hist2)} for the summed bound, with Lemma T and the misprint confirmed and M4 confirmed as reduced; correction: (H-hist2) for $\beta_j$ contains the mixed threshold--history derivative through the running thresholds, so it is a new hypothesis of the (H-hist) kind, not free by a Cauchy estimate in the history.
Open inputs: (H-hist2) (ledger B14, a \status{verified reduction} to the real uniform (H-hist), ledger B13); M3; M4 (units, and the small-field analogue of CMP 122 II (1.73)); the transfer of Ba{\l}aban's printed large-field factor to small-field shell terms; the frozen-integer convention L3; ledger B15.

\paragraph{\file{SOURCES_HHIST2.md}: (H-hist2), second history derivatives.}
For the frozen-threshold map $F=B_j(\cdot;\bar u)$ on $|w-\bar u|\le3r$ in $u$-slots, $\sum_{i,l}|\partial_{u_i}\partial_{u_l}F|\le\varepsilon_j''=C_H(1-\lambda^{1/2})^{-2}u_{j-1}^{-2}$, so $\sum_j\varepsilon_j''\le C_H(1-\lambda^{1/2})^{-2}\gamma^2/\beta$ uniformly in $j$, $K$, and volume, with product decay $\lambda^{m/2}$; the first-derivative half is $\varepsilon_j'\le\tfrac12C_1(1-\lambda)^{-1}u_{j-1}^{-3/2}$, from the real uniform (H-hist). For $\beta_j$ with running thresholds the (H-hist2) step needs only a uniform bound, but Lemma H$'$ still needs a summable first derivative of $\beta_j$ from the real (H-hist).
Route 2(a) differentiates the real chain twice: past slots stay real, the input variable is handled by Cauchy estimates (Theorem SM), and only the own slot is complexified, on one relative $u$-disc (\file{SOURCES_SM.md} \S4.2). The Jacobians $\partial_{u_i}\partial_{u_l}=\tfrac14g_i^3g_l^3\partial_{g_i}\partial_{g_l}$ turn the $1/x$ losses into $g^4$, which is summable because $u_{j-1}\ge\gamma^{-2}+\beta m$.
Verifier: \file{SOURCES_HHIST2_VERIFY.md}, \status{verified reduction} to the real uniform (H-hist), no new input; the Jacobian cancellation and the decay $\lambda^{m/2}$ are confirmed, the product form $\lambda^m$ is \status{refuted} (toy recursion), and $\varepsilon''\propto u^{-2}$ follows from the $u$-disc alone; correction: ``$\beta_j$ needs only a uniform bound'' holds for the (H-hist2) step only, since Lemma H$'$ needs a summable $\beta_j'$.
Open inputs: the real uniform (H-hist) (ledger B13: (Q)/$c_-$, ($M_r$), CMP 102 (189), (M-B) bookkeeping); Theorem SM and the joint $u$-disc and input-ball holomorphy [R]; the large-field part of the frozen map; for $\beta_j$, the mixed threshold--history derivative through Lemma D and M3 (ledger B14).

\paragraph{\file{SOURCES_SM.md}: Theorem SM, the step-map Lipschitz bound.}
The small-field step map factors as $G=\hat G\circ\mathcal L$ with $\mathcal L$ linear, and Theorem SM gives $|DG(I_0)[W]|\le\theta_0\sum_{l\le k}L^{-\alpha_*(k-l)}\|W^{(l)}\|\,e^{-\kappa'd_{k+1}(X)}$ with $\theta_0=2K_HK_V\varepsilon_1$ and $\alpha_*=\min(1,\alpha)$. With it Lemma HH holds at rate $\lambda=L^{-\alpha_*}+C_V\theta_0$ when $C_V\theta_0<1-L^{-\alpha_*}$, and (Eps) of \file{SOURCES_D3W8.md} becomes a corollary with $\varepsilon'=\theta_0/(1-L^{-\alpha_*})$.
The per-input bound CMP 116 (1.24) gives the linear stage a fluctuation factor $\alpha_2^{-1}\varepsilon_1$ and a label weight $L^{-(k-l)}$ (Fact A). The map $\hat G$ is holomorphic on a ball of radius $R_V=\alpha_4\alpha_6e^{-5\kappa}/(4(L+2)^4O(1))$, free of $\varepsilon_1$ and $E_0$, by CMP 116 (C1)--(C5) (Fact B), and a Cauchy estimate on this ball gives the contraction.
Verifier: \file{SOURCES_SM_VERIFY.md}, \status{verified reduction} (items 1--4 confirmed on page images, items 5--6 corrected); the own-slot constant $C_\beta$ has the same $1/g$ gap as $S_0$, so (G-gain) must be applied in both the hop and the own-slot bound, and the per-label decay of the lower-order terms stays reduced.
Open inputs: Fact A for the lower-order bracket terms (ledger A10, now Theorem~\ref{thm:A10} in the corrected form A10$'$; the per-label exponent is $\alpha_*=\beta_H$, a free parameter, and (Eps) reads $\varepsilon'=\theta_0(1-L^{-\beta_H})^{-1}$); Fact B, that is (W4) on CMP 116 (C1)--(C5) with $\alpha_4,\alpha_6$ fixed before $\varepsilon_1$; (I)/L3 (ledger B12); (G-gain) in both slots, carried by \file{SOURCES_GGAIN.md} and \file{SOURCES_GGAIN_LF.md} (ledger B13).

\subsection{Step 7: history bounds and the gain (ledger B13)}\label{sec:src-B13}

\paragraph{\file{SOURCES_HHIST.md}: Lemma HH, (H-hist) in column-sum form.}
Given CMP 109 (0.29), uniform own-slot bounds $\|\partial_{g_{l-1}}V^{(l)}\|_l\le S_0$ and $|\partial\beta_l/\partial g_{l-1}|\le C_\beta$, and a step-map Lipschitz bound (SM) with $\theta<1-L^{-\alpha}$, Lemma HH gives $|\partial\beta_j/\partial g_i|\le(C_\Pi/C_V)S_0\lambda^{j-1-i}$ for $i\le j-2$, with $\lambda=L^{-\alpha}+\theta$. The column sum is then at most $C_\beta+(C_\Pi/C_V)S_0\lambda/(1-\lambda)$, uniformly in $i$, $j$, $K$, and the volume.
The past enters step $k+1$ only through the curly bracket of CMP 109 (2.12) and only through the irrelevant remainders $V^{(l)}$ (the marginal parts cancel, p.~298), and a step-$l$ remainder costs $L^{-\alpha(k-l)}$ per $\pi_k$ cube by (0.29)--(0.30). An induction on $s_l=\|\partial_{g_i}V^{(l)}\|_l$ with one hop per step then gives the rate $\lambda$ in place of the one-hop rate $L^{-\alpha m}$.
Verifier: \file{SOURCES_HHIST_N1_VERIFY.md} (Part 1), \status{verified reduction}: the quotes, the dependence chain, the cube count, and the slot definition are confirmed; the rate is sharp, but the numerical claim ``ratio 1.0'' is false ($s_{i+1+n}=(\theta/\lambda)\lambda^nS_0$, so the bound is loose by $\theta/\lambda$), and Ba{\l}aban gives no per-hop Lipschitz constant $\theta$.
Open inputs: the own-slot bounds $S_0$ and $C_\beta$ ($C_\beta$ is Ba{\l}aban's inductive assumption, ledger A5) and (SM), now carried by \file{SOURCES_GGAIN.md} and \file{SOURCES_SM.md}; validity of the step bounds off the trajectory, close to (I) (ledger B12); the large-field contributions, now in \file{SOURCES_GGAIN_LF.md}; the whole chain is ledger B13.

\paragraph{\file{SOURCES_GGAIN.md}: (G-gain), small field.}
For every step $k$ and $g_k\in\ioc{\gamma}$ the own-slot derivatives obey $\|\partial_{g_k}V^{(k+1)}\|\le C_VS_0$ and $|\partial\beta_{k+1}/\partial g_k|\le C_\beta:=C_\Pi S_0$, with $S_0$ free of $1/g_k$ and of $\varepsilon_1$ (Theorem B2-unif). With Theorem SM the column sum of Lemma HH is uniform, and in $u$-slots it is $O(g_i^3)$, which meets the threshold $Cg_i^{2+\delta}$ of \file{SOURCES_HBETA.md} with $\delta=1$.
By CMP 116 (1.34) the vertex $V_k'(Y,\cdot,B)$ is an analytic function of $x=g_kB$ on a polydisc of $g$-free radius $\varepsilon_1$, so a pointwise Cauchy estimate in $x$ costs $|B|_Y/\varepsilon_1$, the $\varepsilon_1$ cancels against the vertex size, and the one polynomial insertion is absorbed by the Gaussian at strength $\alpha_4$. The $P^{(k)}$ part is cubic in $B'$ by (1.39)--(1.40), so $|\partial_g|\le C_P|B|^3$ there, and in the hop the factor $g_k|B|$ of (1.22)--(1.24) gives $\theta_0\to\theta_gg_k$.
Verifier: \file{SOURCES_GGAIN_VERIFY.md}, \status{proved+verified, conditional on Lemma G3 and the margin ($M_r$)}, for the small-field step on $\ioc{\gamma}$ including vacuum subtraction (items 1, 2, 4, 5, 6 confirmed, 2 and 6 with corrections, item 3 corrected); ($M_r$) is also needed for the $P^{(k)}$ part, Lemma G4 fails as stated and is off-chain, the sharp-cut surface constant is $O(\gamma_2^{-(n_b+1)/2}\varepsilon_1^{-1})$, and Ba{\l}aban's p.~264 assumption is proved only in its first-derivative part, on $\ioc{\gamma}$, for the small-field step.
Open inputs: Lemma G3, hence (Q)/$c_-$ (round 5: proved relative to listed inputs, Theorem~\ref{thm:G3N2}); the margin ($M_r$), which shrinks the $\chi_k$ threshold (ledger C9); the large-field half (\file{SOURCES_GGAIN_LF.md}); the endpoint $g=0$ (ledger C19); all inside ledger B13.

\paragraph{\file{SOURCES_GGAIN_LF.md}: (G-gain), large field.}
With the thresholds frozen at Ba{\l}aban's trajectory $x^*$, the own-slot derivatives of the large-field outputs are uniform in $g_k\in\ioc{\gamma}$ with $|g_k-x^*_k|\le\rho_0x^*_k/2$, in $k$, $K$, and the volume, for example $|\partial_{g_k}R'^{(k)}(X)|\le20\gamma^{2A_0-1}e^{-\kappa d_k(X)}$ and $|\partial_{g_k}T_k(Y_i)1|\le20\gamma^{A_0-1}e^{-\frac12p_0(x^*_k)}$ (Theorem B2-full). With \file{SOURCES_GGAIN.md} and Theorem SM this makes the real uniform (H-hist) a reduction in both halves.
On the disc $|g-x|<\rho_0x$, $\rho_0=0.1$, of \file{SOURCES_HHOL_LF.md} the printed bounds hold with $\gamma_0\to0.826\gamma_0$, and one Cauchy step gives $|\partial_gF|\le20x^{-1}\sup|F|$. The factor $1/x$ is absorbed by the large-field suppression, since $p_0(x)=A_0(\log x^{-2})^{p_0}$ with $p_0$ a positive integer and $A_0$ large (CMP 119 p.~244) gives $x^{-1}e^{-\theta p_0(x)}\le x^{2\theta A_0-1}$.
Verifier: \file{SOURCES_GGAIN_LF_VERIFY.md}, \status{verified reduction} (items 1, 2, 4, 5 confirmed, item 3 corrected); the mismatch average sits at $u'=2x^{-2}-u\in[0.892,1.093]x^{-2}$, which gives the constant $2.48\le2.6$, the Gaussian normalization adds a factor $e^{0.01n}$ per unit cube inside (M-B), and ``$\beta$ from E-class terms only'' is upgraded to a verified reading of CMP 119 p.~279.
Open inputs: (Q)/$c_-$, the margin ($M_r$), CMP 102 (189) through (M-A) (ledger A4), and the (M-B) bookkeeping (ledger B13); not covered: (1.103) at complex $u$, $C^1$ at $g=0$, higher derivatives.

\paragraph{\file{SOURCES_G3.md}: Lemma G3, the $\chi_k$ surface term.}
Given (Q) with constants $c_-$ and $\Lambda$, for every $k$ and $g_k\in\ioc{\gamma}$ the $g_k$-derivative acting on the sharp cut $\chi_{k,Y_0}\chi^c_{k,P}$ of CMP 116 (2.3) in a term of (2.14) is at most $c_1'U(g_k)e^{\delta\kappa d_k(Z_0)}T_{\mathrm{maj}}$, with $U(g)=A_nD_n\varepsilon_1^{-1}e^{-\varepsilon_1^2/(8\Lambda g^2)}$ and $n=N^2-1$. The cut adds $S_\chi=c_2A_nD_n\varepsilon_1^{-1}e^{-\varepsilon_1^2/(8\Lambda\gamma^2)}$ to $S_0$, uniform in $g_k$, $X$, $k$, and the volume, and $O(\gamma^p)$ for every $p$.
Route 2 keeps Ba{\l}aban's cut, a round sphere $|B(b)|=\varepsilon_1/g_k$ per bond, and computes the surface term in polar coordinates, where the Gaussian marginal turns $g^{-1}(\varepsilon_1/g)^n$ into $\varepsilon_1^{-1}t^{n+1}e^{-ct^2}$, $t=\varepsilon_1/g$, with no $1/g$ left. The bond and polymer sums use (2.30), and the smooth cut $\hat\chi$ (Route 3, $\varepsilon_1\to\varepsilon_1/2$, $\gamma\to\gamma/2$) is the fallback.
Verifier: \file{SOURCES_G3_VERIFY.md}, \status{proved+verified, conditional on (Q)} (items 2, 4, 5 confirmed, items 1 and 3 corrected); the surface term is real in the fixed-$B$ computation that \file{SOURCES_GGAIN.md} uses (finite-difference check to 10 digits at $n=8$), (Q) must be read on the joint $(X,B)$ majorant with $a=0$, the bond count is $c_d|Y_0|+|P|$, the printed $\gamma_k$ fails under both parses, and the $B$ versus $B'$ cut question the verifier added is closed in \file{SOURCES_Q.md}.
Open inputs: (Q)/$c_-$ (row B-S.Q; ledger B13), now discharged for the small field by Theorem~\ref{thm:G3N2}; numerical values of $c_-$ and $\Lambda$; higher derivatives.

\paragraph{\file{SOURCES_Q.md}: input (Q), the covariance lower bound $c_->0$.}
The $B(b)$-block of the covariance of Ba{\l}aban's joint Gaussian majorant CMP 116 (2.23) is bounded below by $c_->0$ and above by $\Lambda\le2(1+\gamma_0^{-1})$ for $\alpha_5(1+\gamma_0^{-1})\le\frac12$, uniformly in $k$, the volume, and regular $U$; after the localization step of \file{SOURCES_LOC.md}, $c_-=(\gamma_1+\varepsilon_M)^{-1}$, and in the free case $c_-\ge(a\kappa_dL^{d+1})^{-1}$. The cut of CMP 109 (2.9) is $|B(b)|<\varepsilon_1/g_k$ on every bond except $b_0$, a round sphere with no Jacobian.
The precision of (2.23) is $Q=Q_0-\alpha_5\,\mathrm{diag}(Z,Z_0)\le Q_0$, and inversion reverses order, so the $B(b)$-block of $Q^{-1}$ dominates that of $Q_0^{-1}$, which is $C_0+C_0\Gamma\Gamma^*C_0\ge C_0$ with $C_0=C^{(k)}(Z_0,0)$. The remaining norm input $\|C^*\Delta_kC\|\le\gamma_1$, which CMP 116 (2.7) itself uses, is traced to CMP 99 (3.132), and the free-case constants come from CMP 95 (1.100) and CMP 96 (2.125).
Verifier: \file{SOURCES_Q_VERIFY.md}, \status{verified reduction} to CMP 99 (3.132), (E$_U$), and LOC (items 1, 3, 4 confirmed, 1 and 3 with corrections, item 2 corrected in favour of the reduction); $Q_0$ is coupled, so the block identity is only an inequality, the operator is the $\sigma=0$ localized one (handled in \file{SOURCES_LOC.md}), CMP 99 p.~422 prints the kernel bound (3.132), and $\kappa_d$ depends on $a$ and $L$ unless $a=1$.
Open inputs: the A1 residue (torus-side hypothesis, two conventions, reading $a=1$); the two [I] bridges $Q(U)=Q_{\mathrm{free}}(1+O(\alpha_0))$ and $\|C(U)\|^2\le L^{d+1}(1+O(\alpha_0))$ (ledger A18); (E$_U$), reduced to $k$-uniform kernel moments that Ba{\l}aban does not print; the inputs of LOC. Round 5 assembles (Q) from A7 and A8 instead (Theorem~\ref{thm:G3N2}); its $\Lambda$ corrects 4.3 of this file, because $G_\Gamma$ enters the $B$-block of $Q_0^{-1}$.

\paragraph{\file{SOURCES_LOC.md}: Lemma LOC, the localized covariance $C^{(k)}(Z_0,0)$.}
Given $\gamma_0\le C^*\Delta_kC\le\gamma_1$ and $C_0^{-1}=Z_0C^*\Delta_kCZ_0-R$ with $\|R\|\le\varepsilon_M\le\gamma_0/2$, the localized covariance $C_0=C^{(k)}(Z_0,0)$ of CMP 116 (2.23) has $C_0(b,b)\ge c_-I$ on every bond $b\subset Z_0$ and $\|C_0\|\le2/\gamma_0$. Here $c_-=(\gamma_1+\varepsilon_M)^{-1}$ with the symmetry (S) of the truncated walk sum and $c_-=(\gamma_0-\varepsilon_M)/(\gamma_1+\varepsilon_M)^2$ without it, and $\varepsilon_M=S_d(O(1)e^{-\delta_0M/3}+O(\alpha_0+\alpha_1))$ with $S_d\le d\coth(\delta_0/(4\sqrt d))^d$.
The printed (2.5) is conditioning on $Z_0^c$, so $C^{(k)}(Z_0)$ is the inverse of the compression of $C^*\Delta_kC$ and inherits its bounds, while $C^{(k)}(Z_0,0)$ is the walk sum with the walks meeting $\sigma_0$ deleted and is not a compression. Its distance to the compression is Ba{\l}aban's bound (2.16) applied to $R_2$ at $\sigma\equiv1$, turned into the norm bound $\varepsilon_M$ by the Schur test.
Verifier: \file{SOURCES_LOC_VERIFY.md}, \status{proved+verified, conditional on the $R_2$ bound (2.16), $\gamma_0$, and $\gamma_1$} (item 1 confirmed, items 2--4 corrected); the symmetry (S) is not needed, $\varepsilon_M\to0$ needs $\alpha_0,\alpha_1\to0$ as well as $M\to\infty$, $\gamma_0$ sits in the prose after CMP 99 (3.158) and is proved there only at $U=1$, and the replacement point is the real configuration $U$.
Open inputs: the $R_2$ bound (2.16), which Ba{\l}aban states with ``similarly'' (ledger A2, now a verified reduction to (W$^c$), A7, and two [I] items); the upper bound $\gamma_1$ through CMP 99 (3.132) and (E$_U$); $\gamma_0$ for general $U$ (ledger A3, now a verified full proof, $\gamma_0\ge m_L/(48d^2L^d)$, with residue R-$\Delta$2); $\alpha_0,\alpha_1$ small after $M$ is fixed.

\subsection{Step 8: Ba{\l}aban-level propagator inputs (ledgers A1--A4)}\label{sec:src-A}

\paragraph{\file{SOURCES_A1_KERNEL.md}: ledger A1, kernel-decay form.}
Given (K) and (C) below, the file proves CMP 99 (3.132) in its printed two-scale form, $|(QGQ^*)^{-1}(y,y')|\le 2K_{\mathrm{ext}}L^{d/2+1}(L^j\eta)^{-2}(L^{j'}\eta)^{-d}e^{-\delta_1 d(y,y')}$ for $y\in\Lambda_j$, $y'\in\Lambda_{j'}$, uniformly in $j$, $k$, $\eta$, the volume and regular $U$, and the same for $(QG_1Q^*)^{-1}$ plus a diagonal $a$-term. The rate comes from $\varepsilon=\min\{\delta/2,\,c_-/(2c_+\Sigma(\delta/2))\}$ and is much smaller than Ba{\l}aban's rate. Lemma CT, a Combes--Thomas argument on $\hat N=S^{-1/2}(QGQ^*)S^{-1/2}$ on $\ell^2(\mathfrak B,\mu)$, turns kernel decay (K) of $QGQ^*$ and a lower bound (C) $\hat N\ge c_-$ into decay of the inverse, with no Fourier analysis and no use of A3. (K) is taken from CMP 99 Theorems 3.3 and 3.12 for $G$, (C) from one feasible extension in the variational form $\langle B,N^{-1}B\rangle=\min\{\langle A,G^{-1}A\rangle: QA=B\}$, and the global perturbation from $U=1$ and the Demko--Moss--Smith route are shown not to work. Verifier: \file{SOURCES_A1_A4_VERIFY.md} (Part A), \status{verified reduction}; it confirmed Lemma CT, the two-scale form, the non-circular order of the $G_1$ clause, and both refutations, and it found the file's face-midpoint bumps for (C) not feasible ($Q_j(U)$ is supported on both whole blocks, CMP 98 (14), CMP 99 (3.102)), so its own grade was ``reduced to (C), open''; the newer \file{PLAN.md} entry regrades A1 to a verified reduction once (C) is supplied by \file{SOURCES_QE_INV.md}, \file{SOURCES_QE_KE_HQ.md}, and \file{SOURCES_S2_FORM.md}. Open inputs: the residue of ledger A1, namely the torus-side hypothesis (side $\ge c\cdot RM$ coarsest blocks), two unprinted conventions, and the reading $a=1$ of CMP 99 Theorem 3.1 [I]; the acceptance of the small rate $\delta_1$ by each downstream consumer (S4) is not re-checked (low risk).

\paragraph{\file{SOURCES_QE_KE_HQ.md}: (KE$'$) and (H$_Q$), inputs of (C).}
(KE$'$): Ba{\l}aban's $H$ of CMP 96 Cor.\ 2.8 already satisfies $Q_{98}(1)H=I$ on all of $\mathfrak B$, interfaces and $\Lambda_0$ included, with the decay of CMP 96 (2.151) for $H$ and $\nabla H$, because CMP 96 uses the straight-line average CMP 95 (1.11)/(1.18). (H$_Q$): $|K_j(U;c,b)-q_j(c,b)\operatorname{Ad}U(\Gamma_{c_-\leftarrow b})|\le c_H\alpha_0\ell_j^2\eta_j^d$ on $B^j(c_-)\cup B^j(c_+)$, $c_H=2C'_1+c_T(d)L$, uniformly in $j\le k$. For (KE$'$) the file reads the pages: the invariance group $Q'_j\lambda=0$ of CMP 96 (2.7) fits only the straight-line average, and CMP 98 p.\ 27 states that the propagators hold for the average (1.11). For (H$_Q$) it applies Ba{\l}aban's induction (144)--(145) to the difference between $Q_j(U)$ and the product of the main terms, and it compares the averaged background links with fine straight-line transports and fixed axial contours by a cone-area bound. Verifier: \file{SOURCES_QE_KE_HQ_VERIFY.md}, \status{proved+verified}, (KE$'$) by citation on the page images (CMP 95 (1.11), (1.12), (1.18); CMP 96 (2.7), (2.20), (2.35); CMP 98 p.\ 27 verbatim), so the gap of \file{SOURCES_QE_INV_VERIFY.md} item 2 was a misreading; corrected: the verifier proved the one [R] step (iii) itself (ribbon of $\le(d-1)L^2$ plaquettes, lasso decomposition, error $\le9.6(d-1)\alpha_0\ell_i^2$ for $\alpha_0\le\min(c_2,1/(8(d-1)))$), the cited (112) is the wrong reference and is not needed, and the corridor fallback $H+d\psi$ is exact within layers but not across interfaces (not needed). Open inputs: the unvalued constants $C'_1(d,L)$ and $c_{\mathrm{cone}}(d)$, and two unprinted conventions ($U(c)$ the straight segment; translation-covariant axial contours), both in the residue of ledger A1.

\paragraph{\file{SOURCES_QE_INV.md}: Theorem QE, input (C) of Lemma CT.}
For the free multiscale minimizer dressed with axial transports, $\mathcal E_U$, the file proves, given (KE) and (H$_Q$), $\|Q(U)\mathcal E_U-I\|_S\le c_QM\alpha_0+c_{\mathrm{far}}e^{-\delta_HRM/(16L^2)}$ in the $S$-weighted norm on $\mathfrak B$, uniformly in $j$, $k$, the volume, $\{\Omega_j\}$, and regular $U$. Hence $K_Q\le2$, and $A=\mathcal E_U(Q(U)\mathcal E_U)^{-1}B$ is an exact extension, $Q(U)A=B$ and $A=0$ on $\Lambda_0$, with $\langle A,G^{-1}A\rangle\le K_{\mathrm{ext}}\|B\|_S^2$, $K_{\mathrm{ext}}=4(1+O(M\alpha_0))K_{\mathcal E}+a$, which gives (C) with $c_-=1/K_{\mathrm{ext}}$. The defect of $Q(U)\mathcal E_U$ is written exactly as a $q|E|$-weighted sum of loop holonomies minus one, and a cone-area lemma bounds each holonomy linearly in the distance. A Schur test with the decay (KE), a near/far split across layers with far terms $e^{-cRM}$, and a Neumann series then give $K_Q\le2$, and a crude absolute-value Schur bound gives the energy. Verifier: \file{SOURCES_QE_INV_VERIFY.md}, \status{proved+verified, conditional on (KE$'$), (H$_Q$), S2, and a torus side $\ge c\cdot RM$}; numerics \status{numerical, reproduced} (defect linear in the curvature, slope $0.010$ at $d=2$ and $0.019$ at $d=4$, rigorous Schur constant 41--76 times the true slope); corrected: $\varepsilon''$ is of order $M\alpha_0$, not $\alpha_0$, $\alpha_F=2\cdot O(1)M\alpha_0$ with $O(1)\ge10$ (so the Neumann threshold reads $O(1)M\alpha_0\lesssim0.08$), and a torus-size hypothesis must be added. Open inputs: (KE$'$), (H$_Q$), and S2 for $G$ have since been closed (\file{SOURCES_QE_KE_HQ_VERIFY.md}, \file{SOURCES_S2_FORM_VERIFY.md}); what stays open is the torus-side hypothesis and the two unprinted conventions of ledger A1.

\paragraph{\file{SOURCES_S2_FORM.md}: S2, the form bound for $G$.}
For every bond field $A$, $|\langle A,\Delta'_\pi A\rangle|\le C'_\Pi M\alpha_0(\|DA\|^2+\|D^*A\|^2+\|\ell^{-1}A\|^2)$, hence $G^{-1}\le(1+O(M\alpha_0))(D^*D+DD^*)+O(M\alpha_0)\ell^{-2}+Q^*aQ$ and $\langle\mathcal E_UX,G^{-1}\mathcal E_UX\rangle\le K(d,L,a)\|X\|_S^2$, uniformly in $j$, $\eta$, the volume, $\{\Omega_j\}$, and regular $U$. The route through CMP 99 (3.131) stalls, because (3.131) is printed for one-block supports and a tiling then forces sharp cut-offs that cost $\eta^{-1}$ on the sup side. Route $\Pi$ instead uses the exact identity (3.119), $\langle A,\Delta_\pi A\rangle=\langle\tilde A,\Delta\tilde A\rangle$ with $\tilde A=A-D\lambda$, $\lambda=G'RD^*A$, the curl of a gradient $\eta^{-2}(\operatorname{Ad}U(\partial p)-1)\lambda$, and the block $L^2$ bounds (3.46) of Theorem 3.1 with CMP 96 Lemma 2.1 and a Schur test, which give $\|\ell^{-2}\lambda\|,\|\ell^{-1}D\lambda\|\le B_0c_1\|D^*A\|$. Verifier: \file{SOURCES_S2_FORM_VERIFY.md}, \status{proved+verified, conditional on (KE$'$), $a=1$ in Theorem 3.1, and the block geometry}, together with \S5.1 and \S5.3 of \file{SOURCES_QE_INV.md}; corrected: $\Delta'$ of (3.10) has a second term $\eta^d\tr(DA)^2(\mathrm{Re}\,U(\partial p)-1)$, bounded by $O(M\alpha_0)\|DA\|^2$, so $C'_\Pi=C_\Pi+2c_R(1+\kappa^2a_0^2)$ and $c_-$ drops only by a factor $1-O(M\alpha_0)$; (3.120) carries a minus sign; the normalization of $Q^*aQ$ is settled by (3.16) (one constant $a$). Open inputs: the reading $a=1$ of Theorem 3.1 and the block-geometry item (both in the residue of ledger A1); the $G_1$ analogue, left at sketch level here, is closed in \file{SOURCES_S2_G1_V7.md}.

\paragraph{\file{SOURCES_S2_G1_V7.md}: S2 for $G_1$, and the V7 repair of ledger A2.}
Part A proves $|\langle B,\Delta^{(2)}B\rangle|\le c_2M\alpha_0\|\ell^{-1}B\|^2$, $c_2=\dim\mathfrak g\cdot C_{137}$, and $|\langle A,\Delta^{(2)}_\pi A\rangle|\le2c_2M\alpha_0(\|\ell^{-1}A\|^2+B_0^2c_1^2\|D^*A\|^2)$, so $\langle\mathcal E_UX,G_1^{-1}\mathcal E_UX\rangle\le K_1\|X\|_S^2$ with $K_1=K(1+O(M\alpha_0))$. Part B proves $|R_2(b,b')|\le(O(1)e^{-\delta_0M/3}+O(\alpha_0+\alpha_1))e^{-(1-\theta)\delta_A r}$ for $M\ge\max\{M_*,96\kappa_1/\delta_0\}$ and $\alpha_0+\alpha_1\le\alpha_*$. In Part A, (3.135) is the sandwich $(I-DG'RD^*)^*\Delta^{(2)}(I-DG'RD^*)$, so the form is taken at the $\tilde A$ of (3.119), and the kernel of $\Delta^{(2)}$ is local by (3.137) and symmetric as a Hessian by (3.136), which allows a Schur test with AM--GM. In Part B, $C(\sigma;U_c)=A_Z(U)^{-1}-W'$ with $W'=W_\sigma(U_c)+[C(1;U_c)-C(1;U)]$, the bracket is bounded by a one-variable Cauchy estimate on the complex walk bound, and a Neumann series around the real $A_Z(U)$ lets $A_Z(U)$ cancel in $R_2$. Verifier: \file{SOURCES_S2_G1_V7_VERIFY.md}, Part A \status{proved+verified, conditional} on the conditions of \file{SOURCES_S2_FORM.md}, Part B \status{verified reduction}; corrected: symmetry of $\Delta^{(2)}$ is load-bearing for the column sums, the Cauchy route needs (W$^c$) on the $\alpha'$-space ($\alpha'\gg\alpha$), not only on the $\alpha$-space that $W_\sigma$ needs (already in the ledger, no new item), the two-constant theorem weakens the complex input to a sup bound but does not remove it, touching $\sigma_0$ cubes occur generically, not always, and the two covariances do not coincide on a gauge-fixed subspace; the lower bound for $G_1$ (A3.5) was not checked. Open inputs: for A2, (W$^c$) on the $\alpha'$-space, now proved in the exponential gauge without the Thm 3.15 transfer (\file{SOURCES_WC_EXPGAUGE.md}, \file{SOURCES_WC_RESIDUE.md}, both verified) relative to the inputs listed under \file{SOURCES_A2_R2.md}, the real decay A7 [R], the boundary-contact reading of item (i) [I-forced], and $\operatorname{dist}(Z_0,\Lambda_k^c)>RM$ [I]; for Part A, the inherited reading $a=1$ and the block geometry.

\paragraph{\file{SOURCES_A1_WALK.md}: ledger A1w, walk-expansion form of A1.}
Theorem A1-W: for regular $U$, $(QGQ^*)^{-1}=\sum_\omega T_\omega$ with factors localized in $X\times X$, $U$-dependence on $\tilde X_0^5\cup\dots\cup\tilde X_n^5$, and $|T_\omega(y,y')|\le(4/c_-)L^{d/2+1}(L^j\eta)^{-1-d/2}(L^{j'}\eta)^{-1-d/2}(K_W/M_1)^{|\omega|}e^{-\frac12\varepsilon d(\omega,y,y')}$, absolutely convergent; this is the (3.108) form at rate $\delta_W=\tfrac12\varepsilon\ll\tfrac12\delta_0$, and the same holds for $(QG_1Q^*)^{-1}$ with trees. The local operator on a cube is $G$'s walk sum truncated to $\square^+$, its positivity $\tfrac12c_-$ comes by restriction of the global (C) (not A3), and an accretive Lemma CT gives $\|C_\square\|_\varepsilon\le4/c_-$, as in Ba{\l}aban's own local step (CMP 96 p.\ 237). The parametrix $P=\sum h_\square C_\square h_\square$ satisfies $\hat NP=I+\sum F_\square$ with three remainders (separation, truncation, commutator), each $\le K_F/M_1$, and the Neumann series is the walk series. Verifier: \file{SOURCES_A1_WALK_VERIFY.md}, \status{verified reduction} (construction verified); corrected: (L1) must read ``kernel vanishes within about $M_1$ of $\partial X$'', CMP 96 (2.84) has $e^{-\frac18\delta_0M}$ and (2.82) three error sums, the $\breve G_2$ citation is p.\ 432, the RM condition (2.59) at the walk rate joins the constant order, (W$_G$) has proof sketches printed only for $G'$, $(Q'G'^2Q'^*)^{-1}$ and $G_0$ (Theorems 3.7, 3.9, 3.10), and (C) is conditional, not proved. Open inputs: (W$_G$) [R] (itself reduced in \file{SOURCES_WG.md} to Ba{\l}aban's printed sketches, the Neumann series (3.130), and CMP 96 Lemma 2.1, with WG-1 in the $M_1$ form $\alpha_0\le c'M^{-1}M_1^{-10}$), (C) with the residue of ledger A1, CMP 96 Lemma 2.1, the [I] items (Lipschitz cut-offs O2, block-locality of $Q$, $\kappa_1$ independent of $M$), and O3 ($Z_0$-compressed variant) and O4 (complex $\alpha'$-space), both now closed in \file{SOURCES_WC_EXPGAUGE.md} and \file{SOURCES_WC_RESIDUE.md} (verified), O4 relative to CMP 99 Thm 3.4 [V] and the complex walk form of $G$ [R]. The construction is now written out in full as Theorem~\ref{thm:A1W} (Section~\ref{sec:src-full}); its statement and constants supersede the ones of this paragraph.

\paragraph{\file{SOURCES_WG.md}: (W$_G$), walk expansions of $G$ and $G_1$.}
The file writes out Ba{\l}aban's walk expansions (W$_G$) of the final propagators $G$ and $G_1$ of CMP 99 Sect.~E: for regular $U$, $M\ge M_*(d,L)$ and $M\alpha_0\le a_*(d,L)$, $G=\sum_\omega T_\omega$ is a walk expansion with bounded branching and rate $(1-\theta)\tfrac12\delta_0$, and $G_1$ has a tree expansion with rate $(1-\theta)\tfrac12\varepsilon$, set by the walk expansion of $(QGQ^*)^{-1}$. The literal (3.108) form, with a factor $M^{-1}$ per step also at the $\Delta'_\pi$ junctions, needs in addition $\alpha_0\le c'M^{-11}$ (WG-1; the $M_1$ form is in Theorem~\ref{thm:WG}).
The identity $G^{-1}=G_0^{-1}-\Delta'_\pi$ and the Neumann series (3.130) $G=\sum_nG_0(\Delta'_\pi G_0)^n$ are expanded with the printed walk expansions of $G_0$, $G'$ and $(Q'G'^2Q'^*)^{-1}$ (CMP 99 Thms 3.10, 3.7, 3.9), one local $J$-insertion of size $O(M\alpha_0)$ per $\Delta'_\pi$, and a concatenation lemma that splits the exponent at each junction and sums the junction points by CMP 96 Lemma 2.1. For $G_1$ the same scheme runs on $G_1=\sum_nG(\Delta_\pi^{(2)}G)^n$ with $\Delta_\pi^{(2)}=(I-DRG'D^*)\Delta^{(2)}(I-DG'RD^*)$, whose side branches use the walk expansion A1-W of $(QGQ^*)^{-1}$, and the order $G_0\to G\to(QGQ^*)^{-1}\to H\to G_1$ has no loop.
Verifier: \file{SOURCES_WG_VERIFY.md}, \status{verified reduction}: (W$_G$) for $G$ reduces to the printed proof sketches of CMP 99 Thms 3.7, 3.9, 3.10, the Neumann series, CMP 96 Lemma 2.1 and [I] items, and the $G_1$ clause to the same plus A1-W (Ba{\l}aban himself writes ``assuming (3.132)'' on p.~426); corrections: WG-1 needs $\alpha_0\le c'M^{-11}$, not $M^{-8}$, because a $\Delta'_\pi$ insertion has up to $N_V=9$ walk blocks, and ``after adjusting a definition of $\delta_0$'' is on p.~427.
Open inputs: WG-1 (admissible in Ba{\l}aban's constant order, not implied by $M\alpha_0\le a_0$); A1-W for the $G_1$ clause, with (C) conditional on the torus-side hypothesis and two conventions; the [I] items ($Q$-locality, two-sided derivative bounds, $C^{(2)}$-locality, block counts); all under ledger A1w (O3 and O4 are closed, see \file{SOURCES_A1_WALK.md} above). The write-out \file{WG_FULL.md} (Theorem~\ref{thm:WG}, Section~\ref{sec:src-full}) supersedes the statement and constants of this paragraph: the walks live on the $M_1$-cover, and WG-1 reads $\alpha_0\le c'M^{-1}M_1^{-10}$ for $G$ and $\alpha_0\le c''M^{-1}M_1^{-12}$ for $G_1$.

\paragraph{\file{SOURCES_A2_R2.md}: ledger A2, $R_2$ and $R_3$ of CMP 116 (2.16).}
The precision difference $R_2:=C_0^{-1}-\operatorname{Re}C^{(k)}(Z_0,\sigma)^{-1}$ obeys $|R_2(b,b')|\le(O(1)e^{-\delta_0M/3}+O(\alpha_0+\alpha_1))e^{-\rho|b-b'|}$ for $|\sigma(\Delta)|\le e^{\kappa_1}$, with the printed rate $\rho=\tfrac12\delta_0$ when the decay rate $\delta_A$ of $C^*\Delta_kC$ exceeds $\tfrac12\delta_0$, and a smaller rate $\rho<\delta_A$ otherwise, which (2.21) and LOC tolerate. $R_3$ holds in two weaker forms (R3a), (R3b), sufficient for its only printed use (2.21), and LOC's input (E) follows directly at $\sigma=0$ with $\varepsilon_M=O(e^{-\delta_0M/2})$.
A Neumann series around the compression $C(1)^{-1}=Z_0AZ_0$, $A=C^*\Delta_kC$ (case (a), (2.5)), treats $W_\sigma=C(1)-C(\sigma)$, the walks that meet $\sigma_0$, as the perturbation: each such walk makes a round trip to distance $>\tfrac23M$ from $Z_0$, which gives $O(e^{-\delta_0M/3})$ once the weight $e^{16\kappa_1}$ is absorbed for $M\ge96\kappa_1/\delta_0$. The $O(\alpha_0+\alpha_1)$ part comes from a one-variable Cauchy estimate on the analyticity space, and the verifier's repair V7 expands around the real $A_Z(U)$, so that decay of $A$ is needed at real $U$ only.
Verifier: \file{SOURCES_A2_R2_VERIFY.md}, \status{verified reduction} (with \file{SOURCES_S2_G1_V7.md} and its verifier); corrections: the walk bounds (W) are CMP 99 (3.107)--(3.108), not CMP 109; Lemma N and the Cauchy disc needed decay of $A$ at complex $U$, repaired by V7; the Cauchy radius must respect the plaquette condition; and ``not needed: A1'' is wrong, since (W) for $C^{(k)}(Z_0)$ contains the unprinted walk expansion of $(QGQ^*)^{-1}$ (A1w).
Open inputs: (W$^c$), the walk bounds on the complex $\alpha'$-space, is no longer open: it is proved directly in the exponential gauge, real and complex, with no use of the Thm 3.15 transfer (\file{SOURCES_WC_EXPGAUGE.md}, \file{SOURCES_WC_RESIDUE.md}, both verified; ledger A2 regraded to a verified reduction); it rests on the $G_1$ tree expansion (CMP 99 Thm 3.10, printed sketch), the conditional (C) of A1-W, R2 and R-$\Delta$2 of A3, CMP 109 p.~254 [R], CMP 99 Thm 3.4 [V], and [I] items; it is now written out in full as Theorem~\ref{thm:WC} (Section~\ref{sec:src-full}); the decay of $C^*\Delta_kC$ at real $U$ (A7); the boundary-contact reading of ``intersect'' [I-forced]; $\operatorname{dist}(Z_0,\Lambda_k^c)>RM$ [I]; all under ledger A2.

\paragraph{\file{SOURCES_A3_GAMMA0.md}: ledger A3, the Kato route.}
The lattice Kato (diamagnetic) route to $C^*\Delta_kC\ge\gamma_0$ is refuted, and A3 is reduced instead to Ba{\l}aban's two-step mechanism at $U=1$ made gauge covariant, with $\gamma_0\ge(48d^2)^{-1}L^{-d-1}$ independent of $k$, volume and $M$. This is the verifier's sharpened constant; for $L\ge10$ it carries the factor $Lm_L$ found later in \file{SOURCES_A3_GAMMA0_FULL.md}.
Step (i) bounds the minimal-extension energy below by the unit-lattice curl energy, $\langle B,\Delta_kB\rangle\ge\|\partial_1B\|^2$ at $U=1$ by exact lattice Stokes and Cauchy--Schwarz, and step (ii) is CMP 96 Lemma 2.4, a unit-lattice Poincar\'e inequality in axial gauge. Both steps are made covariant in a local regular gauge on each region, with the gauge-fixing part removed through (3.115) and relative errors $O(M\alpha_0+L^2\alpha_0)c(d,L)$ for $M\ge cL$.
Verifier: \file{SOURCES_A3_GAMMA0_VERIFY.md}, \status{verified reduction} at sketch level, not a proof; corrections: the example $A=d\delta_x$ fails once $\|RD^*A\|^2$ is included and is replaced by an unfixed mode $\lambda\in\operatorname{Ran}(\Delta^{-2}Q'^*)$ ($0$ against $4.62$), so the Kato route stays refuted; $c_*\le1$, not $4$, so the free constant $(12d^2)^{-1}L^{-d-1}$ stands and the bound sharpens from $(192d^2)^{-1}L^{-d-1}$ to $(48d^2)^{-1}L^{-d-1}$; R1 must be CMP 99 Thms 3.1, 3.3, 3.11, 3.12.
Open inputs: R1 with a uniform bound on the oblique projection $\kappa_\pi=\|I-DG'RD^*\|^2$; R2, the $k$-step averaging comparison (an A18 piece); the small items R3--R5; the $\sigma$-localized clause [I]; all under ledger A3, and taken up in \file{SOURCES_A3_GAMMA0_FULL.md} and \file{SOURCES_A3_SIGMA.md}.

\paragraph{\file{SOURCES_A3_GAMMA0_FULL.md}: ledger A3, full proof.}
Theorem A3: for $U$ obeying (3.35)--(3.36) with $M\alpha_0$ small, $\langle\tilde B,C^*\Delta_kC\tilde B\rangle\ge\gamma_0\|\tilde B\|^2$ with $\gamma_0=m_L/(48d^2L^d)$, $m_L=\min(2(1-\cos(\pi/L)),1/L)$, uniform in $k$, volume and regular $U$; this equals $(48d^2)^{-1}L^{-d-1}$ for $L\le9$ and is $0.7555$ times it at $L=13$. The proof step inside CMP 96 (2.126)--(2.127) is false for $L\ge10$ (a lost factor $L$), while the statement of Lemma 2.4 is not refuted and stands with the corrected constant. The failure of (2.127) for every $L\ge10$ is checked in Lean (\file{SlabCounterexample.lean}, Appendix~\ref{app:lean}).
Lemma 5 shows that the $G_1$-minimizer already satisfies $RD^*A_*=0$, because $A\mapsto A-DG'RD^*A$ keeps $QA$ and lowers $\langle A,G_1^{-1}A\rangle$ by exactly $\|RD^*A\|^2$, so the gauge-fixing part drops out and no bound on $\kappa_\pi$ is needed; Lemma 5 is checked in Lean from (F1), (F3), (F4) as hypotheses (\file{GaugeMin.lean}). The two steps of \file{SOURCES_A3_GAMMA0.md} are then written out in full (step (i) with $c_*=1$ and a covariant relative error $2(1+\delta^{-1})\varepsilon_2^2$, step (ii) as CMP 96 Lemma 2.4 reproved with $m_L$ in place of $1/L$ and made covariant in one regular gauge per corridor), and the assembly gives $\gamma_0=\gamma_L/4$ with $\gamma_L=m_L/(12d^2L^d)$.
Verifier: \file{SOURCES_A3_GAMMA0_FULL_VERIFY.md}, \status{proved+verified, conditional on R-$\Delta$2 and R3--R5}, with CMP 99 Thms 3.11--3.12 published and (H$_Q$) \status{proved+verified}; the slab claim against CMP 96 was checked first and confirmed (lowest eigenvalue $L^{-d}\min(2(1-\cos(\pi/L)),1/L)$); corrections: (2.127) fails from $L=10$ with free block bonds and from $L=15$ in axial gauge ($d=2$), not only from $L=30$; fact (E) needs a factor $\sqrt N$; the reason for R-$\Delta$2 must be the form (3.136) plus a block-average bound on $C^{(2)}$.
Open inputs: R-$\Delta$2 [R]; the Schur step of R1 from CMP 99 Thm 3.12, not written out; R3, R4, R5, R-$\Delta'$ and R2$'$; the $\sigma$-localized clause, reduced in \file{SOURCES_A3_SIGMA.md}; all under ledger A3.

\paragraph{\file{SOURCES_A3_SIGMA.md}: ledger A3, $\sigma$-localized clause.}
For real $s\in[0,1]^{\sigma_0}$ the $s$-localized covariance $C_s=C^{(k)}(Z_0,s)$ obeys $\|C_s-C_1\|\le\eta_M=S_dB_We^{-\frac13\delta_0M}$ and, with no symmetry assumption, $\langle v,C_s^{-1}v\rangle\ge(\gamma_0-2\gamma_1^2\eta_M)\|v\|^2\ge\tfrac12\gamma_0\|v\|^2$ for $\eta_M\le\gamma_0/(4\gamma_1^2)$, uniformly in $s$, $Z$, $Z_0$, $k$, volume and real regular $U$, with no $\alpha$-condition. The resolvent bound $\|C_1^{-1}-C_s^{-1}\|\le2\gamma_1^2\eta_M$ also gives LOC's input (E) in norm at real points without (2.16).
The walk sum of CMP 116 p.~3 expands the covariance, $C_s=\sum_\omega s^\omega T_\omega$, so $C_s-C_1=-\sum_{\omega\in W_\sigma}(1-s^\omega)T_\omega$ is dominated termwise by one $s$-free absolute majorant over the walks that meet $\sigma_0$, each of which reaches distance $>\tfrac23M$ from $Z_0$. A Schur bound gives $\eta_M$, and the resolvent identity with $\gamma_0\le Z_0AZ_0\le\gamma_1$ turns it into the lower bound on the precision.
Verifier: \file{SOURCES_A3_SIGMA_VERIFY.md}, \status{verified reduction} to (W) and (N); corrections: (W) is CMP 99 (3.107)--(3.108) with the Cor.~3.8 counting, not CMP 109; the bound without symmetry is the verifier's strengthening of the author's bound $\tfrac23\gamma_0$ under (S$_s$), so (S$_s$) is optional; the uniformity of $B_W$ in $Z_0$ uses part of the exponential decay (constants only).
Open inputs: (W), Ba{\l}aban's assertion (p.~13) that the [13] expansion of $C^{(k)}(Z_0)$ with the large domain $X_0$ obeys (3.108) uniformly in $Z_0$ (A1w enters here [I]; in the exponential gauge this form is Theorem~\ref{thm:WC}(ii), relative to its inputs); the interior reading of ``intersect'' for the zero-length walk, forced by (2.16); (N) $\gamma_0\le A\le\gamma_1$, with $\gamma_0$ from ledger A3 and $\gamma_1$ through A7; complex $\sigma$ stays with A2.

\paragraph{\file{SOURCES_A4_D189.md}: ledger A4, CMP 102 (189).}
CMP 102 (189) bounds the two-point Hessian kernel of the nonlinear action part $V$ at fixed $A'$, and the file derives $|V''(A';b,b')|\le c_V\varepsilon_3e^{-\frac12(1-\alpha)^5\delta_0d(b,b')}$, rate $\ge\tfrac14\delta_0$ for $\alpha\le0.129$, volume-free and on the complex domain (77); (190) follows through the Neumann series (188) at rate $\tfrac18\delta_0$ ($\alpha\le0.159$). Every operator that CMP 102 inverts has the form $I$ plus a small decaying kernel, so a weighted Schur Neumann series (CMP 96 Lemma 2.1) suffices, without Demko--Moss--Smith and without a Hessian lower bound. Lemma D2 shows that the second derivative of the linearizing map solves $(I+\mathfrak N)\mathfrak D_2=C''[(I-H\mathfrak D)\cdot,(I-H\mathfrak D)\cdot]$ with the operator of (68), with two-sided decay at $\tfrac12(1-\alpha)^4\delta_0$, and composition then gives $V''$. Verifier: \file{SOURCES_A1_A4_VERIFY.md} (Part B), \status{verified reduction} to A1 (kernel-decay form), published CMP 102 statements, A18, and one bookkeeping item; corrected: the rates are relabelled ($\tfrac12(1-\alpha)^4\delta_0$ for $\mathfrak D_2$, $\tfrac12(1-\alpha)^5\delta_0$ for $V''$), the name $\mathfrak D_2$ clashes with Ba{\l}aban's (first derivative minus linear term), $V''$ uses (3.132) directly through (87)--(88) and the decay of $\tilde G$ depends on A1 through (3.137) and (3.133), and the CMP 102 rates are fractions of the common $\delta_0\le\delta_1$. Open inputs: A1 (itself a verified reduction), A18, the second-derivative bookkeeping of the $V_0$ terms (90)--(97), and the $\nabla$, H\"older and $\Delta$ versions of (190) [I].

\subsection{Step 8$'$: the three write-outs of the propagator expansions}\label{sec:src-full}

Three files write out in full the constructions that the paragraphs on \file{SOURCES_A1_WALK.md}, \file{SOURCES_WG.md}, and \file{SOURCES_A2_R2.md} reduce to: \file{A1_WALK_FULL.md}, \file{WG_FULL.md}, and \file{WC_FULL.md}. Each was checked by an independent verifier file, fixed, and regraded. Each is paper-grade relative to its tagged inputs. None of them removes an input: they replace ``reduced to a printed sketch'' by a written proof that names its inputs. So the tag stays \status{verified reduction}. The inputs common to all three are the lower bound (C), Ba{\l}aban's printed proof sketches of CMP 99 Thms 3.7, 3.9, and 3.10, and a short list of geometry and convention items [I]. Notation: $M_1$ is the cube size of the walks (CMP 116 p.~3), $M$ is the regularity parameter of CMP 99 (3.36) and the cube size of the $s$-rule, and $R_1M_1\le M$. A walk expansion is a sum $T=\sum_\omega T_\omega$ over walks of domains, with (L1) each factor kernel vanishing within $\tfrac13M_1$ of the boundary of its domain (the band of CMP 116 p.~3), (L2) $U$-dependence of each factor only on $\tilde X^5$, (B1) a per-walk bound $B\rho^{|\omega|}e^{-\delta d(\omega,y,y')}$ with the path distance (3.93), and (B2) bounded branching.

\begin{theorem}[Theorem A1-W, \file{A1_WALK_FULL.md} \S0.5 \status{verified reduction}]\label{thm:A1W}
Let $\hat N=S^{-1/2}QGQ^*S^{-1/2}$. Assume (H1) $\hat N$ has a walk expansion with (L1), (L2), $|\hat n_\omega(y,y')|\le b_G\rho_G^{|\omega|}m(y,y')e^{-\delta_Gd(\omega,y,y')}$, $\rho_G=c_G/M_1$, at most $C_{\mathcal D}$ domains meeting a set of diameter $\le D_+M_1$, and $\operatorname{diam}\tilde X^5\le c_XM_1$; (H2) the lower bound (C), $\operatorname{Re}\langle u,\hat Nu\rangle\ge c_-\|u\|^2$, uniform; (H3) the cube geometry of Ba{\l}aban's cover; (H4) the constant order $(d,L,a,a_0)\to\theta\to c_\pm,\delta_G,\Sigma,b_G,c_G,C_{\mathcal D},c_X\to\varepsilon\to M_1\to M\to\alpha_0$, with WG-1 in the $M_1$ form $\alpha_0\le c'M^{-1}M_1^{-10}$ for $G$ and $\alpha_0\le c''M^{-1}M_1^{-12}$ for $G_1$. Put $c_+=C_{\mathcal D}b_G$ and $\varepsilon=\min\{\tfrac12\delta_G,\ c_-/(8c_+\Sigma(\tfrac12\delta_G))\}$. Then for $M_1\ge M_W(d,L,a,a_0)$: (i) $\hat N^{-1}=\sum_\omega T_\omega$ is a walk expansion; (ii) $|T_\omega(y,y')|\le(4/c_-)(c_1\rho)^{|\omega|}m(y,y')e^{-\frac12\varepsilon d(\omega,y,y')}$ with $\rho\le K_W/M_1$; (iii) $\sum_\omega|T_\omega(y,y')|\le(8C_{\mathcal D}/c_-)m(y,y')e^{-\frac12\varepsilon d(y,y')}$; (iv) $\|\hat N^{-1}\|_{\frac12\varepsilon}\le8/c_-$. All constants are uniform in $U$, $j$, $k$, $\eta$, $\{\Omega_j\}$, and the volume.
\end{theorem}
The file adds two variants. The $Z_0$-compressed variant (\S7): for an operator $A$ with a walk expansion and $\operatorname{Re}A\ge\gamma$, the inverse of $A_Z=1_{Z_0}A1_{Z_0}$ has a walk expansion with prefactor $4/\gamma$, and for $q=M/M_1>10$ and $q\ge3(7+c_X)$ it regroups into the CMP 116 p.~13 form (one domain $X_0$ with $\tilde X_0^5=\tilde Z_0$, all other domains at distance $>\tfrac23M$ from $Z_0$, the band kept, and a merged-slot cost $8C_W/\gamma$ per walk); for $(x+A_Z)^{-1}$, $x\ge0$, the prefactor is $4/(\gamma+x)$. The complex variant (\S8): Lemmas CP and AW and a four-level chain of complex domains give the expansion on the $\alpha'$-space. The chain constants $\Gamma_3=2C_{\mathcal D}B_1'''\Sigma(\delta_1''')$ and $\Gamma_1''=2C_{\mathcal D}B_A''\Sigma(\delta_A'')$ carry the factor $C_{\mathcal D}$, so $k_1$, $k_2$ are smaller by at most $C_{\mathcal D}$. The bound $\sum_\square\varphi_\square^2\ge1$ on $Z_0$, needed in \S7.3, is proved by the verifier (tent family on a grid of spacing $M_1L^k$). The tree form of (H1), which the $G_1$ clause uses, is stated explicitly in \file{A1_WALK_FULL.md} 6.9.
Verifier: \file{A1_WALK_FULL_VERIFY.md}, then a regrade: paper-grade relative to its inputs IN1 ((W$_G$) $=$ Theorem~\ref{thm:WG}), IN3 ((C), conditional), IN8--IN12 [I] (cube geometry, the Lipschitz partition O2, $Q$-locality, the WG-1 condition, $\kappa_1$ independent of $M$), and IN18 [R] ((W$_G$) at complex points, from CMP 99 Thms 3.4 and 3.10).

\begin{theorem}[Theorem WG, \file{WG_FULL.md} 0.7 \status{verified reduction}]\label{thm:WG}
Let $U$ satisfy CMP 99 (3.35)--(3.36) and fix $\theta\in(0,1)$.
(G) (Theorem WG-G.) There are $M_*(d,L,\theta)$ and $a_*(d,L,\theta)>0$ such that for $M_1\ge M_*$, $M\alpha_0\le a_*$, and $RM$ with CMP 96 (2.59) at rate $\theta\tfrac12\delta_0$, $G=\sum_\lambda T_\lambda$ is a walk expansion on the $M_1$-cover, with constants $(B_0',\rho_*,(1-\theta)\tfrac12\delta_0)$, $\rho_*=\max\{c_0M_1^{-1},(K_VM\alpha_0)^{1/10}\}$; the bound holds also with one covariant derivative on either side. If in addition $\alpha_0\le c'M^{-1}M_1^{-10}$ (WG-1), then $\rho_*\le c_0M_1^{-1}$, the CMP 116 (1.7) form.
(G1) (Theorem WG-G1.) Assume in addition Theorem~\ref{thm:A1W} for $(QGQ^*)^{-1}$. There are $M_W$ and $a_{**}>0$ such that for $M_1\ge\max\{M_*',M_W\}$, $M\alpha_0\le a_{**}$, and $RM\ge C(d,L,\theta,\varepsilon)$, $G_1=\sum_\tau T_\tau$ is a tree expansion with rate $(1-\theta)\tfrac12\varepsilon$. The per-step form needs $\alpha_0\le c''M^{-1}M_1^{-12}$.
All constants are uniform in $U$ with (3.35)--(3.36).
\end{theorem}
The proof expands $G=\sum_nG_0(-VG_0)^n$, $V=-\Delta'_\pi$, with the walk expansions of $G_0$, $G'$, and $(Q'G'^2Q'^*)^{-1}$. Lemma N shows that the Neumann series gives the inverse in finite volume, so no separate existence of $G$ is needed. $V$ splits into 24 patterns with at most $N_V=9$ walk blocks, and one local $J$-insertion of size $O(M\alpha_0)$. Lemma C concatenates walks: it splits the exponent at each junction and sums the junction point by CMP 96 Lemma 2.1, so the rate $(1-\theta)\tfrac12\delta_0$ is lost once. The 10 junction steps of a $V$-insertion give the exponent of WG-1. To obtain the band (L1), every domain is enlarged by one layer of $M_1$-cubes; this changes the branching count to $C_{\mathcal D}=5^dc(d)C_{\mathcal D}^{\mathrm{in}}$ and $c_X$ to $c_X+2$, and adds $M_1\ge3c_{\mathrm{loc}}$ to $M_*$. The output satisfies (H1) of Theorem~\ref{thm:A1W}.
Inputs: CMP 99 Thm 3.10 and Thm 3.9 (printed proof sketches), Thm 3.7 with Cor.~3.8 (proof by reference to CMP 95 Prop.~1.2 and CMP 96 Prop.~2.2), CMP 96 Lemma 2.1 under (2.59), the bound $|J|\le O(1)M\alpha_0(L^j\eta)^{-3}$ [V]; for (G1) also Theorem~\ref{thm:A1W}, so (C); [I] items: the scale bookkeeping C0 (written out for the vector $\hat J=S^{3/2}J$), locality of $Q$, $Q'$, $D$, and the parallel transports, the assignment of $\iota$-domains, the locality of $C^{(2)}$, and the $M_1$-cover (asserted on CMP 116 p.~3).
Verifier: \file{WG_FULL_VERIFY.md}, then a regrade: paper-grade relative to these inputs; the verifier found that every consumer uses the $M_1$ form, so the corrected WG-1 replaces the earlier $\alpha_0\le c'M^{-11}$; the band enlargement passed the band regrade.

\begin{theorem}[Theorem WC, \file{WC_FULL.md} 0.10 \status{verified reduction}]\label{thm:WC}
Fix $d\ge2$ and odd $L\ge3$. There are $M_1$, $M$ with $M/M_1>3(6+c_X)$ and $M/M_1>10$, a smallness $\alpha_*$, a four-level chain of complex domains $\alpha'\subset\alpha''\subset\alpha'''\subset$ the $a_1$-domain of CMP 99 Thm 3.4, and constants $\gamma_{\exp}=\tfrac14\gamma_e(d,L)>0$, $\varepsilon''>0$, $\rho>0$, such that for every real regular $U$ and every $Z_0$:
(i) $\langle\tilde B,A_Z\tilde B\rangle\ge\gamma_{\exp}\|\tilde B\|^2$ on the reduced space of $Z_0$;
(ii) $C^{(k)}(Z_0)=\sum_\omega T_\omega$ with (L1), (L2), $|T_\omega(b,b')|\le(4/\gamma_{\exp})(c_1\rho)^{|\omega|}e^{-\frac12\varepsilon d(\omega,b,b')}$, absolutely convergent, and after regrouping in the CMP 116 p.~13 form;
(iii) at every $U_c$ in the $\alpha'$-space, $\operatorname{Re}\langle u,A_Z(U_c)u\rangle\ge\tfrac12\gamma_{\exp}\|u\|^2$, $C^{(k)}(Z_0;\cdot)$ is analytic, and (ii) holds with $4/\gamma_{\exp}\to8/\gamma_{\exp}$ and $\varepsilon\to\varepsilon''$;
(iv) (i)--(iii) hold for $(x+A_Z)^{-1}$, $x\ge0$, uniformly in $x$, with the per-walk prefactor $4/(\gamma_{\exp}+x)$ (real) and $8/(\gamma_{\exp}+x)$ (complex), so the $x^{-1/2}$ integral for $(C^{(k)})^{1/2}$ converges walk by walk.
The constants are uniform in $k$, $\eta$, $Z_0$, $M\ge M_1$, and the background. The proof does not use CMP 99 Thm 3.15 or any transfer from the axial gauge.
\end{theorem}
The lower bound comes from Lemma E, $S_y\le K_b(P_y+G_y)$ on each block, with $K_b=8dL^d(1+16dc_{\mathcal G})\le64d^3L^{2d+1}$ under $16d^3L^dc_{\mathcal G}\tau_p^2\le1$, $\tau_p=16\sqrt N c_4M\alpha_0$; then $\gamma_e^{-1}=K_b(1+6d/m_L)+6/m_L$ and $\gamma_{\exp}=\tfrac12\min(\tfrac12,c_G)\gamma_e$ with $c_G=1$. For Federbush's mean (CMP 109 (0.10)) the exact weights give $c_M\le4\sqrt N$, and $\kappa\le\sqrt2$ for $w\le0.054/\sqrt N$. The walk expansion is the $Z_0$-compressed variant of Theorem~\ref{thm:A1W} applied to the tree expansion of $A_{\exp}$, which rests on Theorem~\ref{thm:WG}(G1). The Combes--Thomas Schur sums carry the factor of $|e^x-1|\le|x|e^{|x|}$: $\Sigma(\tfrac12\delta)\le(1+4/(e\delta))S(\tfrac14\delta)$, which changes only $\varepsilon$, $\varepsilon^a$, and $\varepsilon''$.
Inputs: [R] Theorem~\ref{thm:WG}(G1) (printed sketches), the complex walk form of $G$ (CMP 99 Thms 3.4 and 3.10), the inputs of ledger A3, and CMP 109 p.~254; conditional: (C) and $c_{-,1}\ge c_-(1-O(M\alpha_0))$; [I] items I1--I9 of \file{WC_FULL.md} \S9 (among them the term match of CMP 109 (2.11) with CMP 99 (3.156), the locality of $\tilde C^{(2)}$ and $Q$, the identification $\tilde Q=Q_1(U_1)$, $\operatorname{dist}(Z_0,\Lambda_k^c)>RM$, and the tilde convention).
Verifier: \file{WC_FULL_VERIFY.md} (needs fixes, none structural), then two regrades: after the fixes and after the band fix, paper-grade relative to these inputs.

\subsection{\texorpdfstring{Step 8$''$}{Step 8 bis}: round 5, the remaining Group A rows and Lemma G3}\label{sec:src-r5}

Round 5 treats the Group A rows that had neither a verified proof nor a verified reduction: A7, A8, A10, A11, and A15. It also proves Lemma G3, the residue N2 of the reduction of A5. Each result below has its own project file and an independent verifier file. Each holds relative to the inputs that it names. None of the rows leaves the ledger, because those inputs are not proved here. The tag of each item therefore stays \status{verified reduction} or \status{proved relative to listed inputs}; nothing in this step makes Theorem~\ref{thm:A} or Theorem~\ref{thm:B} unconditional.

\paragraph{The operator of A7.} Ba{\l}aban's $C$ in $C^*\Delta_kC$ is the constraint map $I-hL\tilde Q$, not a square root of the covariance. The operator $C^*\Delta_kC$ is the \emph{precision}: CMP 116 (2.7) reads $(C^{(k)})^{1/2}=(C^*\Delta_kC)^{-1/2}$. So the lower bound $A_Z\ge\tfrac14\gamma_e$ of Theorem~\ref{thm:WC}(i) is the $\gamma_0$ side (row A3) and cannot give $\gamma_1$. The old reduction ``A7 $\Leftarrow$ A1 + A18'' of \file{SOURCES_CDC.md} is superseded; A18 now enters through A3-R4 and the item I3 of \file{WC_FULL.md}.

\begin{theorem}[A7, \file{SOURCES_A7_CDC.md} \status{proved relative to listed inputs; verified}]\label{thm:A7}
Let $U$ be real and regular, $s\equiv1$, and let $\Lambda$ be as in \file{WC_FULL.md} 0.8.
(N) (Lemma N.) $\|C^*\Delta_kC\|\le\gamma_1^N$ with $\gamma_1^N=\|C\|^2\max\{K_{\mathrm{ext},1}-a+c_J+2c_{\mathcal G},\,a+c_J\}$, uniform in $k$, the volume, and $U$.
(D) (Theorem D.) $|(C^*\Delta_kC)(b,b')|\le2B_Ae^{-\delta_A|b-b'|}$ with $\delta_A=(1-\theta)\min\{\tfrac12\varepsilon,\delta_1\}$, and the series (2.7) for $(C^{(k)})^{1/2}$ converges walk by walk for $\gamma_1:=4B_AS(\theta'\delta_A)$.
Both bounds hold for real $s\in[0,1]$ with the same constants.
\end{theorem}
Inputs of (N): C2 of \file{WC_FULL.md} (the bound (C) for $G_1$, conditional on the torus-side hypothesis and two conventions, plus one [I] step for $G_1$), A3-R4, I1, and I3. Inputs of (D): $W_{G_1}$ [R] (Theorem~\ref{thm:WG}(G1)), Theorem~\ref{thm:A1W} [C], I2, and I9. The clause for complex $s$ with $|s|\le e^{\kappa_1}$ and complex $U$ on $\mathcal U(\alpha'')$ is at proof-sketch level and needs $M_1$ large after $\kappa_1$. The value $\gamma_1^N\approx243(K_E+648)$ at $d=4$, $L=3$ is [I].
Verifier: \file{SOURCES_A7_CDC_VERIFY.md}; claims (1), (3), (4), (5) confirmed, claim (2) confirmed with the input list corrected (the two-sided route of \S2.3 is an alternative to Theorem L).

\paragraph{The positivity device of A8.} Only $H(K)=\operatorname{sym}\operatorname{Re}K$ enters the modulus of the Gaussian measure, $K=C^{(k)}(Z_0,\sigma)^{-1}$. Ba{\l}aban states that $C^{(k)}(Z_0,0)$ is symmetric. That claim is not proved: the parametrix walk series is not reversal-symmetric, and the truncation need not keep symmetry. The claim is also not needed. The majorant uses $H_0=\operatorname{sym}C_0^{-1}\ge\tfrac34\gamma_{\exp}$ in its place.

\begin{theorem}[A8, \file{SOURCES_A8_POS.md} \status{proved in corrected form relative to listed inputs; verified}]\label{thm:A8}
(R) (Real case.) At $(U,0)$ with $s\equiv1$, $K=A_Z$ is real symmetric (Lemma S) with $\gamma_{\exp}\le A_Z\le\Gamma_1$, so $d\mu_{C^{(k)}(Z_0)}$ is a positive Gaussian measure.
(D) (Lemma D.) For $|\sigma|\le e^{\kappa_1}$ the $\sigma$-truncation obeys $\|C^{(k)}(Z_0,\sigma)-A_Z^{-1}\|\le\eta_M=(16/\gamma_{\exp})e^{2\cdot3^d\kappa_1}We^{-\varepsilon''M/8}$.
(P) (Lemma P.) The complex Gaussian exists; its modulus is $r_\sigma\,d\mu_{\Gamma_\sigma}$ with $1/(2\Gamma_1'')\le\Gamma_\sigma\le4/\gamma_{\exp}$, which is Ba{\l}aban's (2.15).
(M) (Majorant.) With $H_0=\operatorname{sym}C_0^{-1}\ge\tfrac34\gamma_{\exp}$ the majorant (2.23)--(2.25) holds at a cost $\exp(O(\eta_M^2)|Z_0|)$ and an extra term $O(\eta_M^2)\|ZX\|^2$, absorbed in (2.25).
\end{theorem}
Inputs: Theorem~\ref{thm:WC}(i)--(iii) with their residue (row A2), and five readings I-a--I-e (quadratic forms versus symmetric operators; the $\sigma$-weighted walk sum by the CMP 116 p.~3 rule; three geometry facts; $W<\infty$ at rate $\varepsilon''/8$; ``$\operatorname{Re}C^{-1}$'' as the symmetric real part). The $M$-large conditions are in the order of constants (C24).
Verifier: \file{SOURCES_A8_POS_VERIFY.md}; Lemmas S and P confirmed; Lemma D corrected (a tree-distance slip: the exponent is $e^{-\varepsilon''M/8}$, not $e^{-\varepsilon''M/4}$); the sign of the extra (2.24) term corrected (it is a cost, not a gain). \file{SOURCES_G3.md} reads $R_2$ of (2.16) as $H_0-H_\sigma$.

\paragraph{The form A10$'$ of A10.} Ba{\l}aban writes that the lower-order bracket terms are ``bounded similarly as in (1.24)'' (CMP 116 p.~8). For the raw order-$n$ Taylor terms, $n\le4$, this is false as a bound: the cube count $(6L)^4(L^j\eta)^{-4}$ leaves a sum over $j$ that diverges in $k$. The sentence holds for the terms after the transformations of CMP 109 Sect.~4.

\begin{theorem}[A10$'$, \file{SOURCES_A10_LOWER.md} \status{proved in corrected form relative to listed inputs; verified}]\label{thm:A10}
After the transformations of CMP 109 Sect.~4: the order-0 term is removed by the $t_\square$-derivative; the order-1 term is zero by (4.14); the marginal parts of (4.34) and (5.43) cancel exactly against the counterterm, because $\beta_j$ is defined as the coefficient $\beta$ in (5.43) (p.~298). Every other term obeys the (1.24) form with the factor $8B_0C_1e^{16\kappa_1}\alpha_2^{-1}g_k|B|$ and field power 5, or $4+\beta_H$. The ``other terms'' reach (1.29): F with exponent $(1-2\delta)\kappa$; R and C by the same mechanism.
\end{theorem}
Residue: R1 unprinted constants (finite, fixed before $\varepsilon_1$); R2 the R- and C-terms by mechanism, not written out; R3 Ba{\l}aban's Sect.~4--5 bounds as printed; R4 the tail constant of (4.5) [I]. The per-label exponent of Fact A in Theorem SM is the H\"older exponent $\beta_H$ of CMP 109 (3.31), which Ba{\l}aban does not print (see the paragraph on $\beta_H$ below). Transcription fix: (1.29) reads $O(M^q)$ and $(1/16)\kappa_1$.
Verifier: \file{SOURCES_A10_LOWER_VERIFY.md}; claims (1), (2), (4), (6) confirmed; (3) corrected (a factor 2); (5) corrected ($\beta_H$ is unprinted and needs its own symbol; Lemma HH's $\theta$ has no numerical value, so $\theta<1-L^{-\beta_H}$ holds by $\varepsilon_1$ small).

\paragraph{The free exponent $\beta_H$.} The consumers of A10$'$ change as follows (correction blocks in \file{SOURCES_SM.md} and \file{SOURCES_HHIST.md}): Fact A has $\alpha_*=\beta_H$; (Eps) reads $2\varepsilon'=2\theta_0(1-L^{-\beta_H})^{-1}$; Lemma HH needs $\theta<1-L^{-\beta_H}$, with column sum $\le C_\beta+2C_\Pi\theta_0S_0(1-L^{-\beta_H})^{-1}$; and $\lambda_g=L^{-\beta_H/2}+\theta_g\gamma$ in \S4.3 of \file{SOURCES_SM.md}. The assessment \file{SOURCES_GGAIN_BETAH.md} (\status{unverified assessment}) finds that every consumer needs only $\beta_H>0$, chosen before $\varepsilon_1$ and $\gamma$. Each condition is an $\varepsilon_1$- or $\gamma$-smallness condition. So $\beta_H\in(0,1)$, $\beta_H\le\beta_0<1$, is a free parameter, first in the order of constants (ledger C27). The cost is the factors $(1-L^{-\beta_H})^{-1}$, a larger $B_4(\beta_0)$, and $K_T$ fixed before $\varepsilon_1$. At $L=13$ the (G-gain) margin $1-13^{-\beta_H/2}$ is $0.120$ at $\beta_H=0.1$ and $0.685$ at $\beta_H=0.9$; the Lemma HH margin is $0.901$ at $\beta_H=0.9$. The reading of CMP 102 Thm 1 (9) ``$\beta_0=1$'' as a misprint is [I]. This $\beta_0$ is not the $\beta_0$ of CMP 119 p.~255.

\begin{proposition}[A11, \file{SOURCES_A11_CONST.md} \status{closed as a verified list}]\label{prop:A11}
The ``other constants'' of CMP 119 p.~256 are further scale-indexed numbers $X_j=f(g_j)$: $\alpha_{0,j}$, $\alpha_{1,j}$ ((2.28)), $g_j^{\kappa_0}$ ((2.31)), $\delta_k=(A_1/A_0)\varepsilon_k$ (p.~265), $p_1(g_k)$ ((3.43)), and $g_j(x)$ ((2.24)). Every similar inequality for them follows from (2.6)--(2.9). The three unprinted cross-scale steps ($R_k\le LR_{k+1}$ on p.~268, the monotonicity of $k_1$ on p.~270, and $\varepsilon_{k+1}\le(1+\beta_0)\varepsilon_k$) also follow from (2.9). The side conditions S1--S11 are admissible and acyclic.
\end{proposition}
The file reads 43 introduction points (pp.~255--280 as images, pp.~281--284 as OCR). Three entries join the order of constants: C28 ($M_1<M_2<M$, order $M_1\to M_2\to M$), C29 ($M_2\ge1/\delta$), and C30 ($A_1/A_0\le1/(80B_3)$, $A_1$ after $A_0$). The residue $r>1$ is closed by ``$r\ge2$'' (CMP 119 p.~245). Non-blocking residue: the form of $p_1(g)$ inherited from CMP 122 I, CMP 109 restrictions outside the file, and an OCR slip on p.~278.
Verifier: \file{SOURCES_A11_CONST_VERIFY.md}; claims (1), (4) confirmed; (2), (3) corrected without a change of conclusion (three list omissions, a $(1+\beta_0)^2$ slip fixed by monotonicity, the order $M_1\to M_2\to M$ with $M>M_2$ added).

\paragraph{The whole-lattice coefficient (A15).} $\beta_{j+1}(g_j)$ is defined by CMP 109 (1.20)--(1.22), p.~264: the mixed second moment $\sum_x\Pi_{\mu\nu}(x)x_\mu x_\nu$ of the vacuum-polarization tensor at $U=1$, taken as the limit of periodic-torus kernels $T\nearrow\mathbb Z^d$. Ba{\l}aban states that this limit exists ``by the localized representation (1.7)''. (1.7) and (1.9) say only that each term depends on $U$ restricted to $X$. They do not say that each term is the same function on every torus and on $\mathbb Z^d$. So the existence of the limit stays Ba{\l}aban's bare assertion unless one adds the hypothesis (Loc) below.

\begin{proposition}[Prop.\ A15-N, \file{SOURCES_A15_WHOLE.md} \status{proved given (Loc), (Bd), (H), (Cnt); verified}]\label{prop:A15N}
Fix $j$, $g\in[0,\gamma]$, and $U=1$. Assume (Loc): for every non-wrapping polymer $X$ the torus term equals the $\mathbb Z^d$ term of its lift; (Bd): CMP 109 (1.18); (H): the linearized background map decays exponentially, at the same rate on the torus and on $\mathbb Z^d$; and (Cnt): the standard polymer count. Let $N_j$ be the number of level-$j$ cells per side and $m_j=\tfrac12\min(\delta_H,\kappa-\kappa_0)$. Then the limit in (1.22) exists, and for the torus coefficient $\beta_j^{(N)}$ in moment form
\[ |\beta_j^{(N)}(g)-\beta_j(g)|\le C_jN_j^6e^{-m_jN_j/2}, \]
with $C_j$ independent of $N$. The Fourier form adds $O(N_j^{-2})$.
\end{proposition}
Exact equality $\beta_j^{(N)}=\beta_j$ cannot be derived in general. Image terms and wrapping polymers are real error sources; no explicit $N$ with $\beta_j^{(N)}\ne\beta_j$ was computed, and Ba{\l}aban never claims coefficient equality. Equality holds term by term only for non-wrapping interior $X$ under (Loc). The hypothesis (Loc) is not implied by (L2) of Theorem~\ref{thm:WG}, which covers walk factors only. The case $g=0$ is in \file{SOURCES_T2_G2.md} (ledger B6).
Verifier: \file{SOURCES_A15_WHOLE_VERIFY.md}; (1) confirmed (definition on the page image; the ``(1.41)--(1.43)'' citation on CMP 119 p.~253 is Ba{\l}aban's slip); (2) corrected (three unstated steps: the first-order term drops by gauge invariance, the linearized map includes the $J$ channel, and $m_j$ changes by a fixed factor; the power $N^6$ is loose); (3) corrected (``refuted'' becomes ``cannot be derived''); (4) corrected (the binding-set move below).
Consequence for the binding sets: Theorem 2w uses only the definition of $\beta_j$ (through B5, B6, and A5), so A15 leaves its set. Ba{\l}aban's assertion that the $\mathbb Z^d$ limit of (1.22) exists moves into row A5. A15 stays in the set of Theorem 2w $\Rightarrow$ Theorem 1 and of Theorem~\ref{thm:B}, through B8 and A13: (Loc) for $g>0$, the reading of $\partial\Lambda_{k+1}$ on CMP 119 p.~278, and the R-term part, which cannot hold literally for terms built from merged R$'$ inputs.

\paragraph{The reduction of A5 and Lemma G3.} \file{SOURCES_T2_H0.md} (verified: 47 confirmed, 10 corrected) proves Proposition H0-red: the use of A5 in Theorem 2w, item 7$'$, reduces to a residue N1--N9 in which no constant depends on $j$ or on the past. The residue item N2 is Lemma G3. Ba{\l}aban's $C^\infty$ and all-derivative claims are not used.

\begin{theorem}[Lemma G3 = N2, \file{SOURCES_G3_N2.md} \status{proved relative to listed inputs; verified}]\label{thm:G3N2}
Assume Theorem~\ref{thm:A7} in its two-sided form, Theorem~\ref{thm:A8}, Ba{\l}aban's chain CMP 116 (2.15)--(2.25) read as pointwise exponent bounds (including A2), and the conditions (H$\alpha$) and (H$\gamma$) below. Then (Q) holds in the joint $(X,B)$ form with
\[ c_-=(\gamma_1^N+\varepsilon_0)^{-1}\ge\tfrac{4}{5\gamma_1^N},\qquad \Lambda=2\Big(1+\tfrac{4}{3\gamma_{\exp}}+\tfrac{16}{9\gamma_{\exp}^2}G_\Gamma^2\Big), \]
$\varepsilon_0=2\Gamma_1^2\eta_M^0\le\tfrac14\gamma_{\exp}$, and Lemma G3 holds: the cut adds $S_\chi$ to $S_0$, and $C_\beta=C_\Pi(S_0^{\mathrm{GGAIN}}+S_\chi)$, uniform in $j$, the past, the volume, and regular $U$.
\end{theorem}
The one-bond-fixed form of (Q) is a corollary: fixing $B(b)$ integrates the compression exactly, with Schur complement $(\Gamma_{bb})^{-1}$ and no cost factor ($a=0$). The proof uses no symmetry of $A_Z$ and no Theorem L. It corrects \file{SOURCES_Q.md} 4.3: the $B$-block of $Q_0^{-1}$ is $H_0^{-1}+H_0^{-1}\Gamma\Gamma^*H_0^{-1}$, so $G_\Gamma=\|\Gamma_k(Z_0,0)\|$ enters $\Lambda$.
Conditions. (H$\alpha$): $(\alpha_5+G_\Gamma^2\delta_S)\Lambda_0\le\tfrac12$, with $\delta_S=O(\eta_M^2)$ (ledger C25). (H$\gamma$): $\tfrac1{20}\gamma_2\varepsilon_1^2/\gamma^2\ge4\kappa+1$, and the condition $\tfrac1{20}x-4M^4e^{-x/20}\ge0$ at $x=\gamma_2\varepsilon_1^2/g_k^2$ that the verifier added for (2.31) (ledger C26). Both are smallness conditions of Ba{\l}aban's type, one unit stronger than his own.
The bound on $G_\Gamma$ (\file{SOURCES_G3_GGAMMA.md}, \status{verified}, all four claims confirmed): $\Gamma_k(Z_0,0)=1_{Z_0}A_{\exp}(0)1_{Z_0^c}R(0)$ with $R(0)=(C^{(k)})^{1/2}(0)$, and the Schur test gives $G_\Gamma\le\bar G:=\Gamma_1K_{1/2}$, $\Gamma_1=2B_AS(\delta_A)$, $K_{1/2}=8\gamma_{\exp}^{-1/2}S(\tfrac12\varepsilon)+4\gamma_1^{-1/2}S((1-\theta')\delta_A)$. So $G_\Gamma$ is bounded uniformly in $M$ (verified), and in the volume, $j$, $Z_0$, and regular $U$. With this bound (H$\alpha$) is the single block-size threshold
\[ M\ge\max\{M_\Gamma,M_{\alpha5}\},\quad M_\Gamma=\tfrac{4}{\varepsilon''}\log\big[16h^3\Gamma_1^4\bar G^2\bar\Lambda_0(16W/\gamma_{\exp})^2\big],\quad M_{\alpha5}=\tfrac{3}{\delta_0}\log(8c_5\bar\Lambda_0), \]
$h=4/(3\gamma_{\exp})$, together with Ba{\l}aban's smallness of $\alpha_0,\alpha_1,\alpha_4,\gamma_2$; $\kappa_1$ does not enter $M_\Gamma$, and there is no cycle. One reading residue stays [I]: the $s$-rule acts on Ba{\l}aban's resolvent expansion (3.185) of [13] (CMP 116 p.~13), and the bound identifies $R(0)$ with the walk sum of Theorem~\ref{thm:WC}(iv) with the $\sigma_0$-walks deleted (the reading I-b of A8).
Other inputs of Theorem~\ref{thm:G3N2}: N3 (the Kotecký--Preiss derivative step, the margin ($M_r$) of A17 with C9, dominated differentiation), N6 (Ba{\l}aban's sup induction; past-uniformity of $E_0$ is reduced, not proved), N9 (A15, definition only), and the reading that the pointwise bounds hold on the closed small-field set.
Verifier: \file{SOURCES_G3_N2_VERIFY.md}; (1), (2), (4) confirmed; (3) corrected (the $e^{|P|}$ payment comes from the fifth $e^{-(1/10)x|P|}$ factor via (H$\gamma$), which adds the second $\gamma$ condition); 2000 random numerical instances pass.
Consequence: the use of A5 in Theorem 2w is a \status{verified reduction} with residue N1--N9, and N2 is proved relative to the inputs above. The small-field ``(Q)/$c_-$ open'' of ledger B13 is discharged. A5 stays a row: N5 (continuity at $g=0$, [R]), N6, N8 (readings [I]), and N9 are not proved here, and N7 is B7.

\subsection{\texorpdfstring{Step 8$'''$}{Step 8 ter}: round 6, locality, the boundary, the 2D gap, and the order of constants}\label{sec:src-r6}

Round 6 treats the Ba{\l}aban-level residue of row A15, the R-term row B8, and the open 2D gap of row B4. It also fixes one acyclic order of constants. Each result has its own project file and an independent verifier file. The exception is the check file \file{SOURCES_T2_VOL.md} for (Vol), which has no verifier file. Each result holds relative to the inputs that it names. The tags stay \status{verified reduction} or \status{proved relative to listed inputs}. Nothing in this step makes Theorem~\ref{thm:A} or Theorem~\ref{thm:B} unconditional.

\paragraph{Locality of torus activities.} Work at level $k$ on $T_N=\mathbb Z^4/N\mathbb Z^4$ in level-$k$ units, with projection $\pi$. A set is \emph{liftable} if it lies in $\pi(B)$ for a box $B\subset\mathbb Z^4$ of side $<N/2$. A domain $X$ is \emph{non-wrapping} if its collar $X^\sharp$ ($X$ plus one layer of width $c_k$, the largest cube side at level $k$) is liftable; $\operatorname{diam}_\infty X\le N/2-3c_k$ suffices. $\hat X$, $\hat U$, $\hat J$ denote lifts.

\begin{theorem}[Theorem LOC, \file{SOURCES_A15_LOC.md} \S7.5 \status{proved relative to listed inputs; verified}]\label{thm:LOC}
For every $k$, every coupling history in Ba{\l}aban's range, and every non-wrapping $X\in\mathcal D_{k+1}$,
\[ E^{(k+1)}_{T_N}(X,g,U,J)=E^{(k+1)}_{\mathbb Z^4}(\hat X,g,\hat U,\hat J)\quad\text{for all }(U,J)\in\mathcal U^c_{k+1}(X,\alpha_0,\alpha_1). \]
There is no error term.
\end{theorem}
Method. The mechanism is the $s$-rule of CMP 116 p.~3. At $s=0$ off $Y_0$, every walk-expanded operator is a finite walk sum inside $Y_0$, and the background equations (1.3)--(1.4) decouple by components. A bijection between the surviving torus walks and $\mathbb Z^4$ walks gives equal kernels (Lemma W). Uniqueness gives equal background fields (Lemma N). The cluster assembly and the subtractions are cube geometry (Lemmas A, S). The image terms of a periodized kernel sit in winding walks, so they enter only wrapping activities. The induction on $k$ uses that $\beta_{k+1}$ is the same $\mathbb Z^d$ number in both volumes; a torus coefficient would break exactness at the next step.
Inputs: (FI) each walk factor is given by the same local formula on $T_N$ and $\mathbb Z^4$ (supported by CMP 99 Sect.~C); (U) uniqueness of the solutions of (1.3)--(1.4) (printed in CMP 102); and the Group C rows C32--C36. These are (Cube): $U_j(\square_0,\cdot)$ of [14] uses only data on $\square_0$ [I]; (Cases): the other term classes are localized ``in exactly the same way'' as the typical term (CMP 116 p.~2) [R]; (Conv): cube sides divide $N$, the partitions are $\pi$-images of the $\mathbb Z^4$ partitions, and $M=KR_0M_0$ with $K$ a power of $L$; (Det$'$): a domain-assignment rule for the closed-walk terms of the determinant expansion of CMP 102; and the reading of ``$H_j$ with free boundary conditions'' (CMP 109 p.~290) as the $\mathbb Z^4$ operator [I].
Verifier: \file{SOURCES_A15_LOC_VERIFY.md} (page images); confirmed with minor corrections. (i) Lemma G(b) had the split of CMP 116 (2.7) reversed: for large $x$ the resolvent is a power series in $C^*\Delta_kC$, and the integral with $\tilde G_3(x)$ covers $x\in[0,\gamma_1]$; the conclusion is unchanged. (ii) The determinant expansion of CMP 102 does not follow the $s$-rule. It is a trace of walk expansions, and a trace splits exactly by walk; this needs the convention (Det$'$). (iii) The object table omitted two block-local terms (a Jacobian and $\log\sigma$); they are intrinsic by the same argument. (FI) and (U) are supported on the pages.
Consequences. With Theorem~\ref{thm:LOC}, Proposition~\ref{prop:A15N} holds relative to (Bd), (H), (Cnt), and the inputs above; (Loc) is no longer a hypothesis. Ba{\l}aban's assertion that the $\mathbb Z^d$ limit of CMP 109 (1.22) exists at $g>0$ follows from Proposition~\ref{prop:A15N}. So row A5 has no bare assertion left.

\paragraph{The boundary of the small-field region.}
\begin{proposition}[Proposition BDRY, \file{SOURCES_A15_BDRY.md} \status{proved relative to listed inputs; verified with corrections}]\label{prop:BDRY}
(f) The small-field terms $E_0^{(k+1)}(\Lambda_{k+1},X,z)$ of CMP 119 p.~278 with $X\subset\Lambda_{k+1}$, also those with $X$ next to $\partial\Lambda_{k+1}$, equal the whole-lattice terms exactly. (Thr) Let $N_k$ be the number of level-$k$ cells per side of the torus. If $N_k>22M$, every wrapping or image term of the coefficient analysis at step $k$ lies in Ba{\l}aban's irrelevant class $X\not\subset\tilde\square^2$ (CMP 109 (4.36), CMP 119 (3.48)), and the torus correction to $\beta_j$ is at most $\varepsilon(M)=C_j(22M)^6e^{-11m_jM}$, valid for $m_jM\ge1$.
\end{proposition}
Method. On the image of CMP 119 (1.20), $S_1\cap\Lambda_1=\emptyset$. The change of variables (1.22) acts only on $S_1^*$ bonds, so $H$, the Jacobians, and $V(S)$ reach only terms that meet $\Lambda^c$, and these go to the large-field part. Step $k+1$ (p.~279) keeps every $X\subset\Lambda_{k+1}$. Lemma W$_\Lambda$: with $s=0$ off $Y_0\subset X\subset\Lambda$, every surviving walk of $C(\Lambda)$ lies inside $Y_0$ and does not see $\partial\Lambda$.
Inputs: the inputs of Theorem~\ref{thm:LOC}; C37 (Enl): the $s$-rule localization domain of a walk is $\bigcup X_i^5$ as a whole-lattice set, not $\bigcup X_i^5\cap\Lambda$; (L2$_\Lambda$) through (L2) of Theorem~\ref{thm:A1W}; (Dom); (Conv$_\Lambda$) [I]; and the nesting of the $\Lambda_j$ [I].
Verifier: \file{SOURCES_A15_BDRY_VERIFY.md} (page images). (1) Corrected: Ba{\l}aban prints neither form of the walk domains, so (Enl) is a convention, admissible. Under the $\Lambda$-restricted convention, (f) holds with a margin of at most one $M$-cube layer, also admissible. $C^*\Delta_kC$ is walk-expanded, not finite range (CMP 116 (2.7), CMP 119 p.~275); the $s$-truncation of the $\Delta_k$ walks follows the same convention. (2) Corrected: see the B8 paragraph below. (3) Corrected: the local set is $\tilde\square^5$, of side $11M$, so the threshold is $N_k>22M$, not $10M$. (4) Corrected, admissible: see (Vol) below.

\paragraph{Row B8 replaced.} The old row B8 asked that the finite-volume R-term increment equal the infinite-volume one. That form is ill-posed: the whole lattice of CMP 109 has no R-terms, and no infinite-volume large-field construction exists to compare with. (The author called it false; the verifier corrected this to ill-posed.) CMP 119 p.~279 prints instead that the R$''$ terms with $X\subset\Lambda$ ``do not depend on $\Lambda_{k+1}$''. Claim L$_G$ of \file{SOURCES_A15_BDRY.md} proves this: the step output on $X$ is one map $\mathcal G_X$ of the inputs on $X$, the same in every volume. B8 is replaced by three statements:
(a) $\beta_j$ has no R part, because it is subtracted from $E_{\mathrm{pure}}$ (p.~279, confirmed on the image);
(b) $E_{\mathrm{pure}}$ on a non-wrapping $X\subset\Lambda$ equals the $\mathbb Z^4$ term exactly;
(c) $E_R$ is uniform in $K$ through F2$'$, with the threshold $N_k>22M$.
All three are \status{proved relative to listed inputs; verified}: the inputs of Theorem~\ref{thm:LOC} and Proposition~\ref{prop:BDRY}, the F2$'$ inputs ((W4), (Eps), (D); ledger B20), Ba{\l}aban's (1.18), (3.48), (4.36), and (Vol). B8 is an input of Theorem 2w $\Rightarrow$ Theorem 1, not of Theorem 2w.

\paragraph{The volume hypothesis (Vol).} Proposition~\ref{prop:BDRY} needs $N_k>22M$ at every step $k\le K$. $N_k$ is smallest at $k=K$, where $N_K=\ell$. So the torus must have side $\ell>22M$ final-scale cells. This is (Vol) of Section~\ref{sec:lattice} (ledger C31; \file{THEOREM2_PROOF_v4.md} \S1.1). The verifier graded (Vol) admissible with a real condition. CMP 109 prints the torus exponent (p.~251, $L_\mu=L^m$) and the cube exponent (p.~257, $M=L^m$) with the same letter $m$. If the two were equal, then $\ell=2M$, and (Vol) would fail at the last steps. \file{SOURCES_T2_VOL.md} reads this as reuse of a symbol, and this paper fixes the two exponents independently. (Vol) does not depend on $K$, the proof of Theorem 2w does not use it, and it excludes only tori with $\ell\le22M$.

\paragraph{The 2D gap (G$_A'$).} Let $A'=H'+D'D'^\dagger$ be the charged block of the SU(2) Wilson Hessian at constant flux $f$ on the coarse $\mathbb Z^2$ lattice, in Landau gauge, as in \file{SOURCES_T2_S123.md} \S5. Write $\alpha=ef$, so $|\alpha|=2|f|$.

\begin{theorem}[(G$_A'$), \file{SOURCES_T2_GAPRIME2.md} Theorem 5.1 and Corollary 5.2 \status{proved+verified}]\label{thm:GAprime}
Let $0\le c<1$. For every real $f$ with $0<|f|\le((1-c)/100)^2/2$, the operator $A'$ has no spectrum in $(-2c|f|,2c|f|)$, and neither has $H'$ restricted to $\ker D'^\dagger$. For example, $c=\tfrac12$ gives a gap $|f|$ for $0<|f|\le1.25\cdot10^{-5}$. The constants do not depend on $j$, $L$, or the averaging.
\end{theorem}
Method. A lattice ladder operator $\Pi=\nabla_0+i\operatorname{sgn}(\alpha)\nabla_1$ gives $X(X-2|\alpha|)\ge-157|\alpha|^3$ on low states. So the scalar magnetic Laplacian $\mathfrak L$ has no spectrum in $[(1+\theta)|\alpha|,(3-\theta)|\alpha|]$ for $\theta>26|\alpha|^{1/2}+157|\alpha|$. Then $A'=\mathfrak L\otimes1-2\alpha\sigma_y+R$ with $\|R\psi\|\le48|\alpha|^{3/2}$ near 0, and a contradiction argument gives the gap. The $\ker D'^\dagger$ form follows from $A'D'=D'\mathfrak L$. No Harper-band or Helffer--Sj\"ostrand result is used.
Verifier: \file{SOURCES_T2_GAPRIME2_VERIFY.md}; all four claims confirmed. The constants 26, 157, and 48 are rederived. The lattice commutator is tracked exactly through $2|\sin(\alpha/2)|$. No invariant subspace is needed. The midpoint-transport identification moves from [I] to [P]. The verifier restated the range in units of $f$ (as above) and reproduced the numerics (141 values).

\begin{remark}[the proved range is narrow]\label{rem:GAprime}
The proof covers only $|ef|\le10^{-4}$. Numerics show the gap on all of $\alpha\in(0,3.02]$, with physical ratio $\ge0.95$: at 11 rational fluxes with denominators up to 2213, the smallest $|\text{eigenvalue}|$ of $A'$ is $|\alpha|(1-0.125|\alpha|)$, and a sparse shift-invert run at $\alpha=10^{-4}$ gives the ratio $1.0000125$. So the proof is valid but far from sharp, and the numerics are not a proof. The Nielsen--Olesen mode is the band at $-|ef|(1+0.12|ef|)$, not a zero mode (numerical). The upper edge $3|ef|$ of $H'$ is not proved and is not needed. Open: Prop.~5.6 of \file{SOURCES_T2_S123.md} uses (G$_A'$) for $|f|\le\min(f_0,f_1(j))$, where $f_1(j)$ comes from its size condition $g_f\ge4\eta_1+4C_M^2|f|^3$. No one has checked that this $f$-range fits inside the proved range. The check \file{SOURCES_T2_FRANGE.md} is running.
\end{remark}
Consequence. Open item O1 of \file{SOURCES_T2_S123.md} is closed. Prop.~5.6 and ($Z_j$) hold given the inputs of that file. ($M_j$) (row B4) reduces to (I3) [P], Prop.~4.4 region 1 of \file{BETA_FP_M.md} [argued], and the neutral sector [I]. Both residues are under attack. Row B4 is a verified reduction, subject to the $f$-range check of Remark~\ref{rem:GAprime}.

\paragraph{The canonical order of constants.} \file{SOURCES_ORDER_CANON.md} gives one acyclic order of all constants of the source files. It was verified twice, with amendments E1--E8, and it supersedes the orders of rounds 4 and 5 (ledger C5). \file{SOURCES_ORDER_CHECK.md} found two nominal cycles in the orders as earlier write-outs print them. Cycle A, on the $M_1$ side (\file{WC_FULL.md} 6.9, \file{A1_WALK_FULL.md} 9.1): $M_1\to M\to\alpha'''\to\alpha''\to B_{A''},\delta_{A''}\to M_1$. Cycle B, on the $M$ side: $M\to\alpha'''\to\alpha''\to\alpha'\to\varepsilon''\to M$. Neither cycle is real. A split resolves both: the level constants (rates, prefactors, $\Gamma_n$, $k_n$, $\varepsilon''$) are bounds on domains capped by $a_1$, $a_{**}$, and $M\alpha_0\le a_0$, so they are free of $M$ and of the $\alpha$'s and come before $M_1$; the domain sizes $\alpha'''$, $\alpha''$, $\alpha'$ are smallness choices and come after $M$. The verifier also found forward references and a double count of $M$ in the $J$ scale. The amendments fix them, with two new entries: real $\alpha_0\le\alpha_0'/(CM)$ with $C>c_J$, and $M\ge1/c_J$. The block size $M$ has 32 lower bounds and no upper bound (ledger C24), so the admissible set of $M$ is an up-set. $G_\Gamma$ is bounded uniformly in $M$ (\file{SOURCES_G3_GGAMMA.md} + VERIFY; Section~\ref{sec:src-r5}), so (H$\alpha$) is one more lower bound on $M$ and adds no cycle. One caveat stays: the plaquette row of Lemma M1 is nested by reading only (amendment E6).

\subsection{Step 9: the alternative (E$_U$) route}\label{sec:src-EU}

\paragraph{\file{SOURCES_EU.md}: (E$_U$), uniform block extension (alternative route).}
(E$_U$) asks for a fine field $A$ with $Q(U)A=B$ and covariant energy $\le K\|B\|^2$ for every block field $B$, uniformly in $k$; in the free case $K_0(4)=103.25$, $146.73$, $185.64$, $195.19$ at $N=3,5,13,\infty$ ($=\kappa_4-1$ at $a=1$), and for regular $U$ the dressed extension $A=\mathcal E_U(Q(U)\mathcal E_U)^{-1}B$ gives $K\le K_0(1+c_{2/3}(M\alpha_0)^{2/3}+c_1M\alpha_0)$. The claim that transport and back gives $Q(U)A=B$ exactly is false, and the linear form $K\le K_0(1+cM\alpha_0)$ is not proved.
$K_0$ is computed in momentum space by \file{numerics/k0_d.py} from the CMP 98 (124)--(125) form of $Q(V_0)$, an equal-weight average over $L^d$ straight lines with an $O(L^2\alpha_0)$ remainder, and is checked in real space at $d=2$. For regular $U$ the free extension kernel is dressed by parallel transport, the constraint error $\|Q(U)\mathcal E_U-I\|$ is bounded linearly in the plaquette bound $\alpha_F\le20M\alpha_0$ by lattice Stokes, and IMS patching over flat frames of radius $R\propto\alpha_F^{-1/3}$ gives the power $2/3$.
Verifier: \file{SOURCES_EU_VERIFY.md}, $K_0$ \status{numerical, reproduced}, the bound on $K$ \status{derived sketch}, and (E$_U$) \status{verified reduction} to kernel-moment bounds Ba{\l}aban does not print; corrections: the $N=\infty$ tail is $\propto1/M^3$, not $1/M$; (139) is on journal p.~39; $K_0$ belongs to the straight-line $Q_{98}$, not to the corridor $Q_{95}$; $K\to K_0$ is indicated, not proved; the $k$-step remainder constant is $2C'_1$ by CMP 98 (143).
Open inputs: A21 ($k$-uniform moments and decay of the free kernel $E_0$); the unprinted constant $C'_1$ (A9, a verified reading); the four [I] items (a)--(d) of \file{SOURCES_EU_VERIFY.md} \S5 (uniform moments $\Gamma_j$, $\nu_j$, lattice-sum factors, the boundary cutoff at $\partial\Lambda$, the gauge-fixing term); $Q_{95}$ versus $Q_{98}$; the linear form. The route is an alternative to A1 and is not binding. A21 is now closed, see the next paragraph.

\paragraph{\file{SOURCES_T2_EUMOM.md}: ledger A21, $k$-uniform moments of $E_0$.}
The free minimal-extension kernel $E_0$ of the straight-line average $Q_{98}(1)$ (composition of $k$ one-step averages, $N=L^k$) obeys $|E_0(y,c)|,|\nabla_\nu E_0(y,c)|\le Ce^{-\delta|y-c|}$, with $C$, $\delta$ depending on $d$ only. All moments $\nu_j$, $\gamma_j$, the block sums, and $\Gamma_j$ are bounded uniformly in $k$.
The proof uses symbol identities in momentum space, lower bounds on the symbols, a weighted $L^2$ Combes--Thomas bound uniform in $\eta$, and pointwise regularity from the symbol; it uses no Ba{\l}aban estimate.
Verifier: \file{SOURCES_T2_EUMOM_VERIFY.md}, \status{proved+verified}, 31 claims confirmed, 3 corrected, 1 unsupported (a side remark); it closes open item (i) of \file{SOURCES_EU_VERIFY.md}.
Open inputs: the other [I] items of the (E$_U$) route listed above.

\paragraph{\file{SOURCES_CDC.md}: the norm bound $\gamma_1$ (ledger A7).}
The norm bound $\|C^*\Delta_kC\|\le\gamma_1$, uniform in $k$ and regular $U$, is reduced: by Woodbury $(QG_1Q^*)^{-1}=N_1^{-1}+a$ exactly, so the term $-a\langle B,B\rangle$ of CMP 99 (3.156) cancels and $\gamma_1=L^{d+1}[(1+\varepsilon)K+\varepsilon a+c_J]$ with $\varepsilon,c_J=O(M\alpha_0)c(d,L)$ and $K$ the constant of (E$_U$); in the free case $K\le a(\kappa-1)$, which is $195.19$ at $d=4$, $a=1$. The brief's mechanism, a local difference operator for which unitary $U$ suffices, is \status{refuted}: the leading term $(QG_1Q^*)^{-1}$ is an inverse whose naive coefficients scale as $L^{2k}$, so the regularity (3.35)--(3.36) is needed for the upper bound.
Comparing (3.155) with (3.156) identifies $\langle B,(QG_1Q^*)^{-1}B\rangle$ as the minimum of $\langle A,G_1^{-1}A\rangle$ over $QA=B$, and the split $G_1^{-1}=M_1+aQ^*Q$ of (3.128) gives the exact cancellation of $a$. The remaining minimal energy is bounded by a trial extension of $B$ (the lemma (E$_U$)), the $\Delta'_\pi$ and $J$ parts by $O(M\alpha_0)$ perturbation, and the free constant by the CMP 95 Fourier formula (1.99)--(1.101).
Verifier: \file{SOURCES_CDC_VERIFY.md}, \status{verified reduction}: the Woodbury cancellation is confirmed (the variational form needs only $G_1>0$ and $Q$ onto), but (E$_U$) is not the last input, since CMP 99 (3.132) prints a $k$-uniform kernel bound on $(QG_1Q^*)^{-1}$ that the file omits; corrections: $\|C\|^2\le L^{d+1}$, not $1+L^{d+1}$; $K_0$ is an upper bound only; the free case is proved for single-scale $U=1$ only; $J$ is a second derivative of $A$; the negative-mode branch is dropped, and ``regularity is needed for $\gamma_0$'' is unsupported.
Open inputs: CMP 99 (3.132) (ledger A1), or in its place (E$_U$) (A21, alternative route); A18 for the $J$-term and $\|C(U)\|$; the multiscale correction $e^{-\delta RM}$ to the free constant [I]; the $s$-localized operator $C^{(k)}(Z_0,0)$ that Lemma G3 uses, handled by LOC (A2, A3). Through this file ledger A7 was reduced to A1 + A18. Round 5 supersedes this reduction: Theorem~\ref{thm:A7} proves A7 relative to the inputs of Theorem~\ref{thm:WC}, without the kernel form of A1.
